% This is part of Le Frido
% Copyright (c) 2011-2026
%   Laurent Claessens
% See the file fdl-1.3.txt for copying conditions.


%+++++++++++++++++++++++++++++++++++++++++++++++++++++++++++++++++++++++++++++++++++++++++++++++++++++++++++++++++++++++++++
\section{Les entiers naturels}
\label{SECooPJSYooNYaIaq}
%+++++++++++++++++++++++++++++++++++++++++++++++++++++++++++++++++++++++++++++++++++++++++++++++++++++++++++++++++++++++++++

\begin{normaltext}[\cite{BIBooFVHEooDuxDWi}]	\label{NORMooEntierNaturel}
	Comme le dit la page de Wikipédia consacrée aux nombres entiers, il faut savoir ce dont on parle. Les collégiens français apprennent que les entiers sont les nombres positifs; les Québécois et les Belges y ajoutent les négatifs\dots

	Dans le Frido, on dira \og entier naturel \fg\ voire \og naturel\fg\ pour les positifs, et \og entier relatif \fg\ voire \og entier \fg\ quand on accepte les négatifs. C'est un choix pas standard dont il faut avoir conscience.
\end{normaltext}


%-------------------------------------------------------
\subsection{Importance des naturels dans les ensemble}
%----------------------------------------------------

\begin{normaltext}		\label{NORMooLNLHooLYKwDu}
	Je suis certain qu'on peut définir les naturels comme étant un objet avec une propriété universelle bien particulière dans la catégorie des ensembles\footnote{Évidemment : \url{https://en.wikipedia.org/wiki/Natural_numbers_object}.}.

	En tout cas j'ai l'impression qu'on ne peut rien prouver à propos des ensembles Dedekind-finis\footnote{Un ensemble des Dedekind-infini si il est en bijection avec une partie propre; c'est la définition du Frido de «infini» tout court.} sans passer par une bijection avec une partie des naturels.

	Voici quelques résultats que je doute être possible de démontrer sans passer par les naturels.
	\begin{enumerate}
		\item
		      La proposition \ref{PROPooGJLYooTQtlSV} dit que tout ensemble fini partiellement ordonné possède un élément maximal. Voir aussi \cite{BIBooFUKSooQARznK}.
		\item
		      Si \( E\) est fini, alors son ensemble des parties \( \partP(E)\) est fini.
	\end{enumerate}
\end{normaltext}

%---------------------------------------------------------------------------------------------------------------------------
\subsection{Triplet naturel}
\label{SUBooTripletNaturel}
%---------------------------------------------------------------------------------------------------------------------------

\begin{definition}[\cite{RWWJooJdjxEK}]     \label{DEFooBJBOooWlblAx}
	Un \defe{triplet naturel}{triplet naturel} est un triplet \( (\mN, o, s)\) où \( \mN\) est un ensemble, \( o\) est un élément de \( \mN\) et \( s\) est une application \( s\colon \mN\to \mN\) satisfaisant les propriétés suivantes :
	\begin{enumerate}
		\item
		      \( s\) est injective,
		\item       \label{ITEMooQAKJooGKdJsM}
		      \( s(\mN)=\mN\setminus \{ o \} \)
		\item       \label{ITEMooXPYEooFajywh}
		      Si \( \mA\subset \mN\) est tel que \( o\in \mA\) et \( s(\mA)\subset \mA\), alors \( \mA=\mN\).
	\end{enumerate}
\end{definition}

\begin{normaltext}	\label{NORMooTripletNaturelGratterTheoEns}
	Le théorème suivant est typiquement de ceux qui vont demander de gratter la théorie axiomatique des ensembles avec une certaine précision. Comme nous l'avons dit en introduction \ref{NorooFridoIntro1}, nous n'avons pas l'intention de préciser nos axiomes, mais notre principe de construction est connu et dû à John Von Neumann.

	Nous verrons plus tard, en \ref{NORMooQXASooMXqhjI}, que toute partie infinie d'un triplet naturel fournit un nouveau triplet naturel; il en existe donc plusieurs.
\end{normaltext}

\begin{theorem}     \label{THOooOXMHooXYgMqb}
	Il existe un triplet naturel.
\end{theorem}

\begin{normaltext}[Idée de la construction\cite{BIBooSAYFooEtXedK,BIBooQMEFooEUkxJp}]	\label{NORMooConstructionTripletNaturel}

	Nous effectuons une construction de plusieurs ensembles, puis vérifier que l'ensemble final, muni du nécessaire, forme un triplet naturel.

	\begin{subproof}
		\spitem[Le premier élément]
		%-----------------------------
		On pose \( o = \varnothing \). Dans une axiomatisation raisonnable, l'ensemble vide existe toujours.

		\spitem[Le principe de construction]
		%-----------------------------
		On construit \( A_1 = o \cup \{ o \} \) et on pose \( s(o) = A_1 \).
		Puis on définit \( A_2 = A_1 \cup \{ A_1 \} \), et on pose \( s(A_1) = A_2 \).

		De manière générale\footnote{Ici on voit que nous ne sommes pas en train d'écrire une vraie démonstration parce que «de manière générale» n'est pas quelque chose de précisément défini. En particulier nous ne disposons pas encore de théorèmes formalisant les constructions par récurrence;on serait tenté d'utiliser le théorème \ref{THOooEJPYooZFVnez}.}, une fois un ensemble \( A_i \) construit sur ce principe, on ajoute l'ensemble \( A_i \cup \{ A_i \} \) à notre construction, et on pose \( s(A_i) = A_i \cup \{ A_i \} \).

		Finalement, l'ensemble \( \mN \) est la réunion\footnote{Ici aussi c'est compliqué parce que rien ne garantit qu'il existe un ensemble contenant tous les \( A_i\). Une construction un peu similaire construit les ordinaux, et dans ce cas, il n'existe pas d'ensembles contenant tous les ordinaux.} de tous les ensembles \( A_i \) ainsi construits, et l'application \( s \) est définie par  \( s(A) = A \cup \{ A \} \).

		\spitem[Une remarque intéressante]
		%-----------------------------
		En observant bien, on remarque que \( o \in A_1 \) mais aussi \( o \subset A_1 \). De même, \( A_1 \subset A_2 \) et donc \( o \in A_2 \) tout comme \( A_1 \in A_2 \). On a aussi \( o \subset A_2 \)

		En fait, dans tous les ensembles que l'on construit, on aura toujours un principe d'inclusion, et aussi d'appartenance. Bien que l'intuition nous fait penser que ça fonctionne, la preuve n'est pas évidente.

		\spitem[On a un triplet naturel]
		%-----------------------------
		Montrons que  \( (\mN, o, s)\) est bien un triplet naturel.
		\begin{enumerate}
			\item \( s \) est-elle injective? Si \( A \) et \( B \) sont deux éléments de \( \mN \) avec \( A \neq B \), alors par la remarque ci-dessus, et quitte à inverser \( A \) et \( B \), on a \( A \subset B \). On en déduit que \( B \) est construit après \( A \), et donc \( A \cup \{ A \} \subset B \). Mais, toujours par construction, \( \{ B \} \)  ne peut être dans \( A \), ni être égal au singleton \( \{ A \} \): ainsi \( A \cup \{ A \} \subset B \cup \{ B \} \) mais \( A \cup \{ A \} \neq B \cup \{ B \} \), soit \( s(A) \neq s(B) \).

			\item A-t-on \( s( \mN ) = \mN \setminus \{ o \} \)? Par construction, un élément \( B \in  \mN \setminus \{ o \} \) est de la forme \( A \cup \{ A \} \) pour un certain ensemble \( A \), lui-même dans \( \mN \).
			\item Soit \( \mA \) une partie de \( \mN \) tel que \( o \in \mA \) et \( s(\mA) \subset \mA \). Cette dernière hypothèse nous dit que, dès que \( A \in \mA \), alors \( A \cup \{ A \} \in \mA \). Comme \( \mA \) contient \( o = \emptyset \), il contient aussi \( A_1 \), donc aussi \( A_2 \), et par suite \( \mN \) tout entier. Finalement, \( \mA \subset \mN \) et \( \mN \subset \mA \), donc l'égalité.
		\end{enumerate}
	\end{subproof}
\end{normaltext}

\begin{normaltext}[Définition de \( \eN\)]    \label{NORooDefinitionNChoixTriplet}
	Pour la suite, nous considérons un triplet naturel \( (\mN,o,s)\) et nous notons \( \eN=\mN\). Donc la nature de tous les objets que nous allons considérer à partir de maintenant dépend du choix de triplet naturel que nous faisons à présent. Le théorème \ref{THOooFUXMooJuigHK} nous assurera que peu de choses devraient réellement dépendre de ce choix.

	Nous notons \( 0\) l'élément \( o\) et \( 1\) l'élément \( s(0)\). C'est tout ce dont nous avons besoin dans l'immédiat.
\end{normaltext}

%---------------------------------------------------------------------------------------------------------------------------
\subsection{Applications définies par récurrence}
\label{SUBooApplicationRecurrence}
%---------------------------------------------------------------------------------------------------------------------------

\begin{proposition}[Récurrence\cite{RWWJooJdjxEK}]      \label{PROPooXTRCooKwrWkq}
	Soit un triplet naturel \( (\mN,o,s)\) et une application \( P\colon \mN\to \{ 0,1 \}\) vérifiant\footnote{Les plus \randomGender{pointilleux}{pointilleuses} diront que \( 1\) n'est pas encore défini. Bon j'avoue. Ce qui est important est que \( P\) prenne ses valeurs dans un ensemble contenant deux éléments distincts. Si maintenant vous râlez parce que «deux» est encore moins défini, prenez un ensemble quelconque \( A\) et dites que \( P\) prend ses valeurs dans \( \{ A, \partP(A)\}\). Mais êtes-vous bien \randomGender{certain}{certaine} que \( \partP(A)\neq A\) ?}
	\begin{enumerate}
		\item
		      \( P(o)=1\),
		\item
		      pour tout \( a\in \mN\), si \( P(a)=1\), alors \( P\big( s(a) \big)=1\).
	\end{enumerate}
	Alors \( P(x)=1\) pour tout \( x\in \mN\).
\end{proposition}

\begin{proof}
	Nous posons
	\begin{equation}
		A=\{ x\in\mN\tq P(x)=1 \}.
	\end{equation}
	Cet ensemble vérifie la propriété \ref{DEFooBJBOooWlblAx}\ref{ITEMooXPYEooFajywh}. Donc \( A=\mN\).
\end{proof}

\begin{theorem}[\cite{BIBooZFPUooIiywbk,BIBooMSSFooOOeRKE}]       \label{THOooEJPYooZFVnez}
	Soient \( E\) un ensemble, une application \(g \colon E\to E  \), et \( b\) un élément de \( E\).  Alors il existe une unique application \( f\colon \eN\to E\) telle que :
	\begin{enumerate}
		\item
		      \( f(0)=b\)
		\item
		      \( f\big( s(n) \big)=g\big( f(n) \big)\) pour tout \( n\in \eN\setminus\{ 0 \}\).
	\end{enumerate}
\end{theorem}

\begin{proof}
	Nous commençons par l'unicité. Soient \( f_1\) et \( f_2\) deux telles applications. Nous posons
	\begin{equation}
		A=\{ n\in \eN\tq f_1(n)=f_2(n) \}.
	\end{equation}
	Nous avons \( 0\in A\) parce que \( f_1(0)=f_2(0)=b\).

	Supposons que \( f_1(k)=f_2(k)\). Alors nous avons
	\begin{equation}
		f_1\big( s(k) \big)=g\big( f_1(k) \big)=g\big( f_2(k) \big)=f_2\big( s(k) \big).
	\end{equation}
	Nous en déduisons que \( s(k)\in A\). Autrement dit \( s(A)\subset A\). La définition \ref{DEFooBJBOooWlblAx}\ref{ITEMooXPYEooFajywh} nous indique alors que \( A=\eN\), c'est-à-dire que \( f_1=f_2\).

	Nous montrons à présent l'existence en plusieurs étapes.
	\begin{subproof}
		\spitem[L'ensemble est assez grand]
		Nous considérons l'ensemble \( \mA\) des parties \( A\subset \eN\times E\) telles que
		\begin{enumerate}
			\item
			      \( (0,b)\in A\)
			\item
			      \( (n,x)\in A \Rightarrow \big( s(n),g(x) \big)\in A\).
		\end{enumerate}
		L'ensemble \( \mA\) est non vide parce que \( \eN\times E\in \mA\).
		\spitem[Le plus petit]
		Nous posons
		\begin{equation}
			G=\bigcap_{A\in \mA}A,
		\end{equation}
		et nous prouvons que \( G\in\mA\). D'abord \( (0,b)\in G\) parce que cet élément est dans chacun des \( A\in G\). Ensuite si \( (n,x)\in G\), alors pour tout \( A\in\mA \) nous avons \( (n,x)\in A\) et donc \( \big( s(n),g(x) \big)\in A\). Par conséquent \( \big( s(n),g(x) \big)\in \bigcup_{A\in\mA}A=G\).

		Pour \( n\in \eN\) nous posons
		\begin{equation}
			G_n=\{ x\in E\tq (n,x)\in G \}.
		\end{equation}
		Nous avons en particulier que \( b\in G_0\) parce que \( (0,b)\in G\).
		\spitem[\( G\) contient un \( (n,x)\) pour tout \( n\)]
		Nous prouvons que pour tout \( n\in \eN\), il existe \( x\in E\) tel que \( (n,x)\in G\). Nous faisons ça avec la proposition \ref{PROPooXTRCooKwrWkq} en posant
		\begin{equation}
			\begin{aligned}
				P\colon \eN & \to \{ 0,1 \}                           \\
				n           & \mapsto \begin{cases}
					                      1 & \text{si } G_n\neq\emptyset \\
					                      0 & \text{sinon.}
				                      \end{cases}
			\end{aligned}
		\end{equation}
		Puisque \( (0,b)\in G\) nous avons \( P(0)=1\). Supposons que \( P(k)=1\) et montrons que \( P\big( s(k) \big)=1\). Comme \( P(k)=1\), il existe \( x\in E\) tel que \( (k,x)\in G\). De ce fait, \( \big( s(k),g(x) \big)\in G\), ce qui donne \( G_{s(k)}\neq \emptyset\) et \( P\big( s(k) \big)=1\).
		\spitem[\( G_n\) est un singleton]
		Nous avons vu que \( G_n\) n'est jamais vide. Nous allons montrer que \( G_n\) est un singleton pour tout \( n\). Pour cela nous posons
		\begin{equation}
			\begin{aligned}
				P\colon \eN & \to \{ 0,1 \}                                        \\
				n           & \mapsto \begin{cases}
					                      1 & \text{si }  G_n \text{ est un singleton} \\
					                      0 & \text{sinon. }
				                      \end{cases}
			\end{aligned}
		\end{equation}
		Nous prouvons par récurrence que \( P(n)=1\) pour tout \( n\).
		\begin{subproof}
			\spitem[\( P(0)=1\)]
			Nous commençons par prouver que \( P(0)=1\). Nous savons que \( (0,b)\in G_0\). Supposons \( a\neq b\) tel que \( (0,a)\in G_0\). Alors en posant \( G'=G\setminus\{ (0,a) \}\) nous avons \( G'\in \mA\).

			En effet \( (0,b)\in G'\) parce que \( (0,b)\in G\) et \( (0,b)\neq (0,a)\). De plus si \( (n,x)\in G'\), alors \( \big( s(n),g(x) \big)\in G\). Mais comme \( s(n)\neq 0\) nous avons \( \big( s(n),g(x) \big)\neq (0,a)\) et donc \( \big( s(n),g(x) \big)\in G'\).

			L'ensemble \( G'\) serait un élément de \( \mA\) strictement inclus dans \( G\). Impossible. Donc \( G_0\) est un singleton.

			\spitem[Récurrence]
			Supposons que \( P(k)=1\), c'est-à-dire que \( G_k\) est un singleton. Soit \( e\) l'unique élément de \( G_k\) : \( (k,e)\in G\). Nous avons alors aussi que \( \big( s(k),g(e) \big)\in G\). Nous devons prouver que si \( y\in G_{s(k)}\), alors \( y=g(e)\).

			Supposons donc \( y\neq g(e)\) soit dans \( G_{s(k)}\). Nous posons
			\begin{equation}
				G'=G\setminus\{ \big( s(k),y \big) \}.
			\end{equation}
			Nous prouvons que \( G'\in\mA\). D'abord \( (0,b)\in G'\) parce que \( s(k)\neq 0\). Soit ensuite \( (m,z)\in G'\). Si \( m=k\), alors \( z=e\) (parce que par hypothèse \( G_k\) est un singleton) et nous savons que \( \big( s(m),g(e) \big)\in G'\). Si par contre \( m\neq k\), comme  \( s\) est injective, nous avons aussi \( s(m)\neq s(k)\). Donc \( \big( s(m),g(z) \big)\neq \big( s(k),y \big)\) et \( \big( s(m),g(z) \big)\in G'\). Donc \( G'\in \mA\) et est strictement plus petit que \( G\). Contradiction.

			Nous concluons que \( G_{s(k)}\) est un singleton, c'est-à-dire que \( P\big( s(k) \big)=1\).
			\spitem[Conclusion]
			Nous avons prouvé que \( G_n\) est un singleton pour tout \( n\).
		\end{subproof}

		\spitem[Et enfin]

		Nous définissons \( f(n)\) comme étant l'unique élément de \( G_n\). Puisque \( (0,b)\in G\) nous avons \( G_0=\{ b \}\) et donc \( g(0)=b\).

		Par définition de \( f\), nous avons \( \big( n,f(n) \big)\in G\). Parce que \( G\in\mA\) nous avons alors
		\begin{equation}
			\big( s(n),g\big( f(n) \big) \big)\in G.
		\end{equation}
		Autrement dit, \( G_{s(n)}= \{  g\big( f(n) \big) \}\). Cela montre que
		\begin{equation}
			f\big( s(n) \big)=g\big( f(n) \big),
		\end{equation}
		et donc que \( f\) vérifie les propriétés demandées.
	\end{subproof}
\end{proof}

\begin{normaltext}
	Le théorème \ref{THOooEJPYooZFVnez} permet de faire des démonstration par récurrence en posant \(P \colon \eN\to \{ 0,1 \}  \) où \( P(n)=0\) si la proposition à prouver est fausse pour \( n\) et \( P(n)=1\) si elle est vraie. Et on démontre que \( P(n)=1\) pour tout \( n\).

	Parfois, implicitement, nous dirons que \( P(n)=1\) si la propriété est vraie pour tous les ensembles de cardinal plus petit ou égal à \( n\).
\end{normaltext}

\begin{normaltext}		\label{NORMooPUOAooGnFjaX}

	Lorsque nous disons faire une récurrence sur le cardinal d'un ensemble, nous posons la fonction \(P \colon \eN\to \{ 0,1 \}  \) définie par \( P(n)=1\) si la propriété est vraie pour tout ensemble de cardinal \( n\) et \( P(n)=0\) si il existe un ensemble de cardinal \( n\) ne vérifiant pas la propriété.

\end{normaltext}

\begin{normaltext}[\cite{MonCerveau}]      \label{NORMooLNXMooWIblPf}
	Le théorème \ref{THOooEJPYooZFVnez} nous permet de définir une suite \( (x_n)_{n\in \eN}\) par récurrence en posant par exemple
	\begin{subequations}
		\begin{numcases}{}
			x_1=b\\
			x_{n+1}=f(x_n).
		\end{numcases}
	\end{subequations}

	Parfois nous pouvons avoir envie de définir une suite par récurrence de telle sorte que chaque nouveau terme dépende de tous les précédents :
	\begin{subequations}
		\begin{numcases}{}
			x_1=b\\
			x_{n+1}=g(x_1,\ldots,x_{n-1}).
		\end{numcases}
	\end{subequations}
	Pour ce faire,  il faut un peu adapter.

	Posons \( E'=E\cup (E\times E)\cup (E\times E\times E)\cup\ldots\). Nous supposons avoir une application \(g \colon E'\to E'  \) qui ne prend ses valeurs que dans \( E\) : \( g(E')\subset E\). De plus nous considérons l'application
	\begin{equation}
		\begin{aligned}
			a\colon \eN & \to E'                                \\
			n           & \mapsto \big( f(0),\ldots,f(n) \big).
		\end{aligned}
	\end{equation}

	Le théorème \ref{THOooEJPYooZFVnez} appliqué à \( g\circ a\) nous donne l'existence et l'unicité d'une application \(f \colon \eN\to E'  \) telle que
	\begin{subequations}
		\begin{numcases}{}
			f(0)=b\\
			f(n+1)=g\big( f(0),\ldots,f(n) \big).
		\end{numcases}
	\end{subequations}
\end{normaltext}


\begin{corollary}[\cite{BIBooZFPUooIiywbk}]       \label{CORooVNHKooRkKtXf}
	Soient deux ensembles \( X,Y\), une application \( \alpha\colon X\to Y\) et une application \( \beta\colon Y\to Y\). Alors il existe une unique application \( H\colon X\times \eN\to Y\) telle que
	\begin{enumerate}
		\item
		      \( H(x , 0)   = \alpha(x)\)   pour tout élément \( x\in X\);
		\item
		      \( H(x , n+1) = \beta( H( x , n) )\) pour tout élément \( x\in X\) et pour tout \( n\in \eN\).
	\end{enumerate}
\end{corollary}

\begin{proof}
	Pour faire le lien avec les notations du théorème \ref{THOooEJPYooZFVnez}, nous notons \( E=\Fun(X,Y)\), \( b=\alpha\in E\) et
	\begin{equation}
		\begin{aligned}
			g\colon E & \to E                 \\
			s         & \mapsto \beta\circ s.
		\end{aligned}
	\end{equation}
	Le théorème \ref{THOooEJPYooZFVnez} donne alors l'existence d'une application \( f\colon \eN\to E\) telle que
	\begin{enumerate}
		\item
		      \( f(0)=b\)
		\item
		      \( f(n+1)=g\big( f(n) \big)\).
	\end{enumerate}
	Nous définissons alors
	\begin{equation}
		\begin{aligned}
			H\colon X\times \eN & \to Y          \\
			(x,n)               & \mapsto f(n)x,
		\end{aligned}
	\end{equation}
	et nous vérifions qu'elle satisfait aux exigences.

	\begin{enumerate}
		\item
		      D'abord nous avons
		      \begin{equation}
			      H(x,0)=f(0)x=b(x)=\alpha(x).
		      \end{equation}
		\item
		      Ensuite, pour \( x\in X\) et \( n\in \eN\) nous avons :
		      \begin{subequations}
			      \begin{align}
				      H(x,n+1) & =f(x+1)x                      \\
				               & =g\big( f(n) \big)x           \\
				               & =g\big( f(n) \big)x           \\
				               & =\big( \beta\circ f(n) \big)x \\
				               & =\beta\big( f(n)x \big)       \\
				               & =\beta\big( H(x,n) \big).
			      \end{align}
		      \end{subequations}
	\end{enumerate}
	Et voilà.
\end{proof}

%---------------------------------------------------------------------------------------------------------------------------
\subsection{Addition sur les naturels}
%---------------------------------------------------------------------------------------------------------------------------

\begin{propositionDef}[\cite{RWWJooJdjxEK,MonCerveau}]      \label{PROPooVFOXooXmwpFh}
	Il existe une unique fonction \( f\colon \eN\times \eN\to \eN\) vérifiant
	\begin{enumerate}
		\item       \label{ITEMooILZSooNYIkYR}
		      \( f(a,0)=a\) pour tout \( a\in \eN\)
		\item       \label{ITEMooZWHQooBAjZyE}
		      \( f\big( a,s(b) \big)=s\big( f(a,b) \big)\) pour tout \( a,b\in \eN\).
	\end{enumerate}
	Pour \( a,b\in \eN\) nous notons \( f(a,b)=a+b\).
\end{propositionDef}

\begin{proof}
	Soit \( a\in \eN\). Le théorème \ref{THOooEJPYooZFVnez} dit qu'il existe une unique application \(f_a \colon \eN\to \eN  \) telle que pour tout \( x\in \eN\),
	\begin{subequations}		\label{SUBEQSooWUGPooSxezFi}
		\begin{numcases}{}
			f_a(0)=a		\label{EQooVOCZooOCHZLn}\\
			f_a\big( s(x) \big)=s\big( f_a(x) \big).	\label{EQooMEYWooEsWofl}
		\end{numcases}
	\end{subequations}
	Montrons qu'une telle application vérifie \( f_0(a)=a\) pour tout \( a\in \eN\). Pour cela nous utilisons la récurrence de la proposition \ref{PROPooXTRCooKwrWkq} en posant
	\begin{equation}
		\begin{aligned}
			P\colon \eN & \to \{ 0,1 \}                   \\
			n           & \mapsto \begin{cases}
				                      1 & \text{si } f_0(n)=n \\
				                      0 & \text{sinon. }
			                      \end{cases}
		\end{aligned}
	\end{equation}
	Nous avons \( f_0(0)=0\) en posant \( a=0\) dans \eqref{EQooVOCZooOCHZLn}. Supposons ensuite que \( P(a)=1\). Nous avons
	\begin{subequations}
		\begin{align}
			f_0\big( s(a) \big) & =s\big( f_0(a) \big)                      \\
			                    & =s(a)                & \text{pcq }P(a)=1.
		\end{align}
	\end{subequations}
	Donc \( P\big( s(a) \big)=s(a)\). Nous en déduisons que \( P(x)=1\) pour tout \( x\in \eN\) et donc que
	\begin{equation}		\label{EQooIKGVooHQMQqX}
		f_0(x)=x
	\end{equation}
	pour tout \( x\).


	Avec ça, nous prouvons ce qu'il faut.
	\begin{subproof}
		\spitem[Existence]
		%-----------------------------------------------------------
		Nous posons \( f(a,b)=f_a(b)\), et nous prouvons les deux conditions.
		\begin{subproof}
			\spitem[Condition \ref{ITEMooILZSooNYIkYR}]
			%-----------------------------------------------------------
			Nous avons \( f(a,0)=f_a(0)=a \) par	\eqref{EQooIKGVooHQMQqX}.
			\spitem[Condition \ref{ITEMooZWHQooBAjZyE}]
			%-----------------------------------------------------------
			Nous avons :
			\begin{equation}
				f\big( a,s(b) \big)=f_a\big( s(b) \big)=s\big( f_a(b) \big)=s\big( f(a,b) \big).
			\end{equation}
		\end{subproof}

		\spitem[Unicité]
		%-----------------------------------------------------------
		Soient \( f\) et \( g\) deux applications vérifiant les conditions. Nous posons
		\begin{equation}
			\begin{aligned}
				f_a\colon \eN & \to \eN        \\
				x             & \mapsto f(a,x)
			\end{aligned}
		\end{equation}
		et \(g_a \colon \eN\to \eN  \) de la même manière. Nous prouvons que \( f_a\) et \( g_a\) vérifient tous deux les conditions \eqref{SUBEQSooWUGPooSxezFi}. Cela donnera \( f_a=g_a\) et donc \( f=g\).

		D'abord
		\begin{equation}
			f_a\big( s(b) \big)=f\big( a,s(b) \big)=s\big( f(a,b) \big)=s\big( f_a(b) \big).
		\end{equation}
		Ensuite nous prouvons par récurrence que \( f_a(0)=a\) pour tout \( a\in \eN\). D'abord \( f(0,0)=0\), ensuite
		\begin{equation}
			f\big( 0,s(a) \big)=s\big( f(0,a) \big)=s(a).
		\end{equation}
		De même \( g_a\) vérifie les conditions.
	\end{subproof}
\end{proof}

\begin{lemma}[\cite{RWWJooJdjxEK}]      \label{LEMooMJMTooOtUuJT}
	Pour tout \( a\in \eN\) nous avons \( s(a)=a+1\).
\end{lemma}

\begin{proof}
	Nous avons :
	\begin{subequations}
		\begin{align}
			s(a) & =s(a+0)        \label{SUBEQooMNBLooTOruhE} \\
			     & =a+s(0)        \label{SUBEQooAGUSooGijYGj} \\
			     & =a+1.          \label{SUBEQooUZQDooWtNBHO}
		\end{align}
	\end{subequations}
	Justifications.
	\begin{enumerate}
		\item
		      Pour \eqref{SUBEQooMNBLooTOruhE} c'est dans la définition \ref{PROPooVFOXooXmwpFh}\ref{ITEMooILZSooNYIkYR} de la somme.
		\item
		      Pour \eqref{SUBEQooAGUSooGijYGj}, c'est dans la définition \ref{PROPooVFOXooXmwpFh}\ref{ITEMooZWHQooBAjZyE} de la somme.
		\item
		      Pour \eqref{SUBEQooUZQDooWtNBHO}. Le symbole «\( 1\)» désigne l'élément \( s(0)\) dans \( \eN\).
	\end{enumerate}
\end{proof}

\begin{proposition}[\cite{RWWJooJdjxEK}]     \label{PROPooTLTSooGNMTmV}
	En ce qui concerne la somme dans \( \eN\).
	\begin{enumerate}
		\item       \label{ITEMooIFFPooXfftfG}
		      La somme est associative et commutative.
		\item       \label{ITEMooSGRVooPAVFYK}
		      L'élément \( 0\) est neutre.
		\item       \label{ITEMooNUTHooJWWzGv}
		      Tous les éléments de \( \eN\) sont réguliers\footnote{Élément régulier pour une opération, définition \ref{DEFooIJIEooZaAdSs}.} par rapport à l'addition, c'est-à-dire que si \( a+x=a+y\) alors \( x=y\).
	\end{enumerate}
\end{proposition}

\begin{proof}
	En plusieurs parties.
	\begin{subproof}
		\spitem[Associative]
		Nous devons prouver que \( (a+b)+c=a+(b+c)\) pour tout \( a,b,c\in \eN\). Pour ce faire, nous fixons \( a,b\in\eN\) et nous prouvons l'égalité demandée par récurrence sur \( c\).

		Pour \( c=0\), nous avons \( (a+b)+c=a+b\) et \( a+(b+c)=a+b\). Donc nous sommes d'accord\footnote{Notez que nous n'avons pas utilisé le fait que \( 0\) était neutre des deux côtés -- chose que nous n'avons pas encore démontré. Nous avons seulement utilisé \( a+0=a\), qui est dans la définition de la somme.}.

		Nous vérifions avec \( s(c)\) :
		\begin{subequations}
			\begin{align}
				(a+b)+s(c) & =s\big( (a+b)+c \big)                              \\
				           & =s\big( a+(b+c) \big)  \label{SUBEQooDDUBooAjuuZq} \\
				           & =a+s(b+c)                                          \\
				           & =a+\big( b+s(c) \big).
			\end{align}
		\end{subequations}
		Justifications.
		\begin{itemize}
			\item Pour \eqref{SUBEQooDDUBooAjuuZq}. C'est l'hypothèse de récurrence. À ce stade, je vous conseille d'être capable de rédiger complètement la récurrence et l'appel au théorème \ref{THOooEJPYooZFVnez}.
		\end{itemize}
		\spitem[Neutre]
		La définition de l'addition contient déjà \( a+0=a\). Nous prouvons par récurrence que \( 0+a=a\) pour tout \( a\in \eN\).

		Pour \( a=0\), l'égalité demandé est correcte : \( 0+0=0\) parce que pour tout \( x\) dans \( \eN\), \( 0+x=x\). Pour \( s(a)\) nous avons
		\begin{equation}
			0+s(a)=s(0+a)=s(a).
		\end{equation}
		La dernière égalité est l'hypothèse de récurrence.

		Nous posons
		\begin{equation}
			A=\{ a\in \eN\tq 0+a=a \}.
		\end{equation}
		Nous avons prouvé que \( 0\in A\) et que \( s(A)\subset A\). Le théorème \ref{THOooEJPYooZFVnez} nous assure alors que \( A=\eN\).
		\spitem[Commutativité]
		Nous fixons \( a\in \eN\) et nous prouvons par récurrence sur \( b\) que \( a+b=b+a\) pour tout \( b\in \eN\). Cela va être décomposé en plusieurs étapes.
		\spitem[\( a+0=0+a\)]
		Pour \( b=0\) c'est correct, car \( b+0=0+b=b\) parce que \( 0\) est neutre.
		\spitem[\( a+1=1+a\)]
		Nous démontrons par récurrence sur \( a\) que \( a+1=1+a\). Avec \( a=0\) c'est déjà fait. Pour les autres,
		\begin{subequations}
			\begin{align}
				s(a)+1 & =(a+1)+1 & \text{lemme \ref{LEMooMJMTooOtUuJT}} \\
				       & =(1+a)+1 & \text{hypothèse récurrence}          \\
				       & =1+(a+1) & \text{associativité}                 \\
				       & =1+s(a).
			\end{align}
		\end{subequations}
		\spitem[\( a+b=b+a\)]
		Nous y voici. Nous fixons \( a\) et nous prouvons par récurrence que \( a+b=b+a\). Pour \( b=0\) c'est déjà fait. Pour les autres,
		\begin{subequations}
			\begin{align}
				a+s(b) & =a+(b+1)                                    \\
				       & =(a+b)+1                                    \\
				       & =(b+a)+1 & \text{hypothèse récurrence}      \\
				       & =b+(a+1)                                    \\
				       & =b+(1+a) & \text{commutativité avec \( 1\)} \\
				       & =(b+1)+a & \text{associativité}             \\
				       & =s(b)+a.
			\end{align}
		\end{subequations}
		Récurrence terminée.
		\spitem[Régularité]
		Nous devons prouver que, pour tout \( a,x,y\in \eN\), si \( a+x=a+y\) alors \( x=y\). Nous allons procéder par récurrence en posant
		\begin{equation}
			A=\{ a\in \eN\tq\forall x,y\in \eN, a+x=a+y\Rightarrow x=y\}.
		\end{equation}
		Puisque \( 0+x=x\) et \( 0+y=y\), nous avons \( 0\in A\). Supposons à présent que \( a\in A\) et montrons que \( s(a)\in A\). Soient \( x,y\in \eN\) tels que \( a+x=a+y\). Nous avons :
		\begin{subequations}
			\begin{align}
				            & \quad s(a)+x=s(a)+y                                \\
				\Rightarrow & \quad s(a+x)=s(a+y)    \label{SUBEQooNJHZooSIKPxN} \\
				\Rightarrow & \quad a+x=a+y          \label{SUBEQooWDJLooJhzIFe} \\
				\Rightarrow & \quad x=y              \label{SUBEQooTLYZooFsMaJD}
			\end{align}
		\end{subequations}
		Justifications.
		\begin{itemize}
			\item Pour \eqref{SUBEQooNJHZooSIKPxN}. En utilisant la définition de l'addition et la commutativité, nous avons \( s(a)+x=s(a+x)\).
			\item Pour \eqref{SUBEQooWDJLooJhzIFe}. Parce que \( s\) est injective; c'est dans la définition \ref{DEFooBJBOooWlblAx} d'un triplet naturel.
			\item Pour \eqref{SUBEQooTLYZooFsMaJD}. Parce que \( a\in A\).
		\end{itemize}
		Nous avons prouvé que \( s(A)\subset A\), et donc que \( A=\eN\).
	\end{subproof}
\end{proof}

\begin{lemma}       \label{LEMooCOMSooEWrumL}
	Nous avons \( 0\neq 1\).
\end{lemma}

\begin{proof}
	Par définition \( 1=s(0)\). Comme \( s\) est à valeurs dans \( \eN\setminus\{ 0 \}\), nous ne pouvons pas avoir \( s(0)=0\).
\end{proof}

\begin{lemma}[\cite{MonCerveau}]       \label{LEMooQBHFooCuCusQ}
	Si \( a+b=0\), alors \( a=b=0\).
\end{lemma}

\begin{proof}
	Soient \( a,b\in \eN\) tels que \( a+b=0\), et supposons que \( b\neq 0\). Par la définition \ref{DEFooBJBOooWlblAx}\ref{ITEMooQAKJooGKdJsM}, nous avons \( b=s(c)\) pour un certain \( c\in \eN\).

	Dans ce cas nous avons \( a+b=a+s(c)=s(a+c)\neq 0\) parce que l'image de \( s\) ne contient pas \( 0\). Hélas, par hypothèse nous avons \( a+b=0\). Nous avons obtenu une contradiction, et nous déduisons que \( b=0\).

	Maintenant que nous savons que \( b=0\), il reste \( 0=a+b=a+0=a\).
\end{proof}

%---------------------------------------------------------------------------------------------------------------------------
\subsection{Ordre sur les naturels}
\label{SUBooOrdreNaturels}
%---------------------------------------------------------------------------------------------------------------------------

\begin{definition}[Relation d'ordre\cite{RWWJooJdjxEK}]     \label{DEFooAXZSooTEMjlV}
	Pour \( a,b\in \eN\), nous notons \( a\leq b\) si il existe \( x\in \eN\) tel que \( a+x=b\).

	Nous notons également \( a<b \) si \( a\leq b\) et \( a\neq b\).
\end{definition}

\begin{propositionDef}[\cite{MonCerveau}]	\label{LEMooVOTAooALatDu}
	Si \( a\leq b\), il existe un unique \( x\in \eN\) tel que \( a+x=b\).

	Nous notons \( b-a\) l'unique élément de \( \eN\) tel que
	\begin{equation}
		a+(b-a)=b.
	\end{equation}
\end{propositionDef}

\begin{proof}
	L'existence est juste la définition \ref{DEFooAXZSooTEMjlV}. En ce qui concerne l'unicité, si \( a+x=b=a+y\), alors la proposition \ref{PROPooTLTSooGNMTmV}\ref{ITEMooNUTHooJWWzGv} donne \( x=y\).
\end{proof}

\begin{lemma}       \label{LEMooWMYPooLTMyWb}
	Nous avons \( a\leq s(a)\) pour tout \( a\in \eN\).
\end{lemma}

\begin{proof}
	Cela est une conséquence du lemme \ref{LEMooMJMTooOtUuJT} : \( s(a)=a+1\).
\end{proof}

\begin{proposition}     \label{PROPooVXBBooZcghrA}
	La relation \( \leq\) est une relation d'ordre compatible avec l'addition.
\end{proposition}

\begin{proof}
	Plusieurs choses à vérifier.
	\begin{subproof}
		\spitem[Réflexive]
		Nous avons \( a\leq a\) parce que \( a+0=a\).
		\spitem[Antisymétrique]
		Soient \( a,b\in \eN\) tels que \( a\leq b\) et \( b\leq a\). Il existe \( x,y\in \eN\) tels que
		\begin{subequations}
			\begin{align}
				b & =a+x  \\
				a & =b+y.
			\end{align}
		\end{subequations}
		En substituant la seconde équation dans la première, \( b=(b+y)+x\) que nous récrivons, en utilisant l'associativité\footnote{Proposition \ref{PROPooTLTSooGNMTmV}\ref{ITEMooIFFPooXfftfG}.},
		\begin{equation}
			0+b=b+(x+y).
		\end{equation}
		En utilisant la régularité, \( 0=x+y\) et donc \( x=y=0\) par le lemme \ref{LEMooQBHFooCuCusQ}. Cela donne alors \( a=b\).
		\spitem[Transitive]
		Si \( a\leq b\) et \( b\leq c\), nous avons \( n,p\in \eN\) tels que \( b=a+n\) et \( c=b+p\). Donc
		\begin{equation}
			c=(a+n)+p=a+(n+p),
		\end{equation}
		ce qui signifie que \( a\leq c\). Notez l'utilisation de l'associativité de la somme, démontrée en la proposition \ref{PROPooTLTSooGNMTmV}\ref{ITEMooIFFPooXfftfG}.
		\spitem[Compatibilité]
		Soient \( a,b,n\in \eN\) tels que \( a\leq b\). Nous devons montrer que \( a+n\leq b+n\). Puisque \( a\leq b\), il existe \( x\in \eN\) tel que \( b=a+x\). Par conséquent,
		\begin{equation}
			b+n=a+x+n=(a+n)+x,
		\end{equation}
		qui signifie bien que \( a+n\leq b+n\).
	\end{subproof}
\end{proof}

\begin{lemma}       \label{LEMooPVRQooXPMKTt}
	À propos d'ordre et de stricte inégalité.
	\begin{enumerate}
		\item       \label{ITEMooGWWFooYGPCZw}
		      Si \( x\leq a\) et \( b\neq 0\), alors \( x<a+b\).
		\item       \label{ITEMooRWGWooAfkrri}
		      Si \( x\leq a\), alors \( x<s(a)\).
		\item       \label{ITEMooWCOIooMWrCag}
		      Si \( x<a\), alors \( s(x)\leq a\).
	\end{enumerate}
\end{lemma}

\begin{proof}
	En plusieurs parties.
	\begin{subproof}
		\spitem[Pour \ref{ITEMooGWWFooYGPCZw}]
		Si \( x\leq a\), il existe \( d\in \eN\) tel que \( x+d=a\). Nous avons alors aussi
		\begin{equation}
			x+d+b=a+b,
		\end{equation}
		ce qui signifie que \( x\leq a+b\). Mais si \( x\) était égal à \( a+b\), nous aurions \( d+b=0\), ce qui impliquerait\footnote{Par le lemme \ref{LEMooQBHFooCuCusQ}.} \( d=b=0\), alors que l'hypothèse stipule que \( b\neq 0\). Donc \( x\neq a+b\).
		\spitem[Pour \ref{ITEMooRWGWooAfkrri}]
		Il s'agit seulement d'utiliser le point \ref{ITEMooGWWFooYGPCZw} avec \( b=1\) et le fait que \( s(a)=a+1\) par le lemme \ref{LEMooMJMTooOtUuJT}.
		\spitem[Pour \ref{ITEMooWCOIooMWrCag}]
		Par hypothèse, il existe \( b\neq 0\) tel que \( x+b=a\). Puisque \( b\neq 0\), il existe \( c\in \eN\) tel que \( b=s(c)\) et donc, tel que
		\begin{equation}
			x+s(c)=a.
		\end{equation}
		En utilisant le fait que \( s(c)=c+1\) ainsi que l'associativité et la commutativité de l'addition (proposition \ref{PROPooTLTSooGNMTmV}\ref{ITEMooIFFPooXfftfG}) nous avons
		\begin{equation}
			a=x+s(c)=s(x)+c,
		\end{equation}
		ce qui prouve que \( s(x)\leq a\).
	\end{subproof}
\end{proof}

\begin{proposition}[\cite{MonCerveau}]	\label{PROPooAKLVooBkRIAT}
	Soient \( a,b,k\in \eN\). Nous avons
	\begin{enumerate}
		\item			\label{ITEMooFMQCooSjuCDV}
		      \begin{equation}
			      (a-k)+b=(a+b)-k.
		      \end{equation}
		\item		\label{ITEMooKHLDooEEXjmF}
		      \begin{equation}
			      (a-k)+k=(a+k)-k=a.
		      \end{equation}
	\end{enumerate}
\end{proposition}

\begin{proof}
	En plusieurs parties.
	\begin{subproof}
		\spitem[Pour \ref{ITEMooFMQCooSjuCDV}]
		%-----------------------------------------------------------

		Nous posons \( A=(a-k)+b\) et \( B=(a+b)-k\). En utilisant l'associativité et la symétrie de l'addition,
		\begin{subequations}
			\begin{align}
				A+k & =(a-k)+b+k             \\
				    & =\big( (a-k)+k \big)+n \\
				    & =a+b
			\end{align}
		\end{subequations}
		parce que par définition \( (a-k)+k=a\). Nous avons aussi
		\begin{equation}
			\big( (a+b)-k \big)+k=a+b.
		\end{equation}
		Donc \( A+k=B+k\) et donc \( A=B\) par \ref{PROPooTLTSooGNMTmV}\ref{ITEMooNUTHooJWWzGv}, nous avons \( A=B\).

		\spitem[Pour \ref{ITEMooKHLDooEEXjmF}]
		%-----------------------------------------------------------
		En utilisant le point \ref{ITEMooFMQCooSjuCDV} avec \( b=k\) nous avons déjà \( (a-k)+k=(a+k)-k \). Or par définition l'élément \( (a+k)-k\) est le \( x\) tel que \( x+k=a+k\). Cela est forcément \( x=a\) par la proposition \ref{PROPooTLTSooGNMTmV}\ref{ITEMooNUTHooJWWzGv}.
	\end{subproof}
\end{proof}

\begin{proposition}[\cite{MonCerveau}]	\label{PROPooJSZHooVDzBDE}
	Plusieurs choses.
	\begin{enumerate}
		\item		\label{ITEMooVSLRooGhGFgX}
		      Soit \( a\leq b\) dans \( \eN\). Si \( k\leq a\) et \( k\leq b\) alors
		      \begin{equation}
			      a-k\leq b-k.
		      \end{equation}
		\item		\label{ITEMooYGOAooWRbKdj}
		      Nous avons
		      \begin{equation}
			      a-k<a
		      \end{equation}
		      pour tout \( k\neq 0\) dans \( \eN\).
	\end{enumerate}
\end{proposition}

\begin{proof}
	Nous commençons par prouver \ref{ITEMooVSLRooGhGFgX}. Vu que \( a\leq b\), il existe un \( y\in \eN\) tel que \( a+y=b\), et ce \( y\) est nommé \( b-a\) (par définition \ref{LEMooVOTAooALatDu} de la soustraction). Voyons maintenant
	\begin{subequations}
		\begin{align}
			(a-k)+y & =(a+y)-k & \text{prop. \ref{PROPooAKLVooBkRIAT}} \\
			        & =b-k     & \text{pcq. \( a+y=b\)}.
		\end{align}
	\end{subequations}
	Cela signifie que \( a-k\leq b-k\).

	Pour \ref{ITEMooYGOAooWRbKdj}. D'abord nous avons \((a-k)+k=a \) par la proposition \ref{PROPooAKLVooBkRIAT}\ref{ITEMooKHLDooEEXjmF}. Cela prouve que \( a-1\leq a\). Ensuite nous supposons que \( a-k=a\) et nous prouvons que \( k=0\). Pour cela nous ajourons \( a\) des deux côtés :
	\begin{equation}
		(a-k)+k=a+k.
	\end{equation}
	En utilisant à nouveau le fait que \( (a-k)+k=a\) nous avons \( a=a+k\). Et nous en déduisons \( k=0\) avec la proposition \ref{PROPooTLTSooGNMTmV}\ref{ITEMooNUTHooJWWzGv}.
\end{proof}

\begin{proposition}[\cite{MonCerveau}]	\label{PROPooHRCOooBkEvud}
	Soient \( a,x\in \eN\). Nous avons \( x\leq a-1\) si et seulement si \( x<a\).
\end{proposition}

\begin{proof}
	En deux sens.
	\begin{subproof}
		\spitem[\( \Rightarrow\)]
		%-----------------------------------------------------------
		Nous supposons que \( x\leq a-1\). Grâce à la proposition \ref{PROPooJSZHooVDzBDE}\ref{ITEMooYGOAooWRbKdj} nous avons \( x\leq a-1<a\).
		\spitem[\( \Leftarrow\)]
		%-----------------------------------------------------------
		Nous supposons que \( x<a\). Par le lemme \ref{LEMooPVRQooXPMKTt}\ref{ITEMooWCOIooMWrCag}, nous avons \( x+a\leq a\). La proposition \ref{PROPooJSZHooVDzBDE} donne alors
		\begin{equation}
			(x+1)-1\leq a-1.
		\end{equation}
		Et maintenant la proposition \ref{PROPooAKLVooBkRIAT}\ref{ITEMooKHLDooEEXjmF} donne \( x\leq a-1\).
	\end{subproof}
\end{proof}

\begin{lemma}       \label{LEMooCSIXooHeuWEd}
	L'élément \( 0\) est l'unique plus petit élément de \( \eN\).
\end{lemma}

\begin{proof}
	Puisque \( 0\) est neutre pour l'addition\footnote{Proposition \ref{PROPooTLTSooGNMTmV}\ref{ITEMooSGRVooPAVFYK}.}, nous avons \( a+0=a\) pour tout \( a\in \eN\) et donc \( 0\leq a\) pour tout \( a\). Cela veut dire que \( 0\) est plus petit que tout élément de \( \eN\).

	En ce qui concerne l'unicité, soit \( z\in \eN\) tel que \( z\leq a\) pour tout \( a\in \eN\). Si \( z\neq 0\), il existe \( x\in \eN\) tel que \( z=s(x)\). Nous avons donc \( z\geq x\) en même temps que \(x\leq z\). Cela implique \( z=x\) (parce qu'une relation d'ordre est symétrique) et donc \( z=z+1\). En utilisant la régularité de \( z\) pour l'addition nous en déduisons que \( 0=1\), ce qui est impossible par le lemme \ref{LEMooCOMSooEWrumL}.
\end{proof}

\begin{lemma}       \label{LEMooJRZKooOMhOkH}
	Si \( a\leq b\) et \( b\leq a\), alors \( a=b\).
\end{lemma}

\begin{proof}
	L'inégalité \( a\leq b\) dit qu'il existe \( x\in \eN\) tel que \( a+x=b\). En mettant cela dans l'inégalité \( b\leq a\) nous trouvons \( a+x\leq a\) qui donne, via la proposition \ref{PROPooVXBBooZcghrA} : \( x\leq 0\). Nous en déduisons que \( x=0\) parce que zéro est l'unique minimum de \( \eN\) par le lemme \ref{LEMooCSIXooHeuWEd}.
\end{proof}

\begin{proposition}     \label{PROPooGCCRooFBYrlo}
	Le couple \( (\eN,\leq)\) est totalement ordonné\footnote{Définition \ref{DEFooVGYQooUhUZGr}.}.
\end{proposition}

\begin{proof}
	Soit \( a\in \eN\). Nous devons prouver que pour tout \( x\in \eN\) nous avons \( x\leq a\) ou \( a\leq x\) (non exclusifs). Nous posons
	\begin{subequations}
		\begin{align}
			A & =\{ x\in \eN\tq x\leq a \}  \\
			B & =\{ x\in \eN\tq a\leq x \},
		\end{align}
	\end{subequations}
	et nous prouvons que \( A\cup B=\eN\) en montrant que \( 0\in A\cup B\) et que \( s(A\cup B)\subset A\cup B\).

	Nous avons \( 0\in A\subset A\cup B\) par le lemme \ref{LEMooCSIXooHeuWEd}.

	Pour étudier \( s(A\cup B)\), nous considérons \( x\in A\cup B\) et nous subdivisons en deux cas selon que \( x\in A\) ou \( x\in B\).

	\begin{subproof}
		\spitem[Si \( x\in B\)]
		Si \( x\in B\), alors \( a\leq x\leq s(x)\) parce que \( x\leq s(x)\) par le lemme \ref{LEMooWMYPooLTMyWb}. Donc \( s(x)\in B\subset A\cup B\).
		\spitem[Si \( x\in A\)]
		Si \( x\in A\), il y a deux possibilités : \( x=a\) et \( x\neq a\). Si \( x=a\), alors \( a\leq s(x)\) et donc \( s(x)\in A\subset A\cup B\).

		Si \( x\neq a\), alors le lemme \ref{LEMooPVRQooXPMKTt}\ref{ITEMooWCOIooMWrCag} nous indique que \( s(x)\leq a\) et donc \( s(x)\in A\subset A\cup B\).
	\end{subproof}
	Nous avons donc prouvé que \( s(A\cup B)\subset A\cup B\), et donc que \( A\cup B=\eN\).
\end{proof}

\begin{proposition}[\cite{RWWJooJdjxEK,MonCerveau}]     \label{PROPooMZOWooHmsXzI}
	L'ensemble ordonné \( (\eN,\leq)\) vérifie les propriétés suivantes.
	\begin{enumerate}
		\item       \label{ITEMooJLAHooDKukfH}
		      L'élément \( 0\) est l'unique minimum de \( \eN\).
		\item       \label{ITEMooYAJIooEFmOpB}
		      Toute partie non vide a un unique plus petit élément.
		\item       \label{ITEMooSRGOooNYJJHY}
		      L'ensemble \( \eN\) n'a pas de plus grand élément.
		\item       \label{ITEMooKIHZooDRTCdx}
		      Toute partie non vide majorée a un unique plus grand élément.
	\end{enumerate}
\end{proposition}

\begin{proof}
	Point par point.
	\begin{subproof}
		\spitem[Pour \ref{ITEMooJLAHooDKukfH}]
		C'est le lemme \ref{LEMooCSIXooHeuWEd}.

		\spitem[Pour \ref{ITEMooYAJIooEFmOpB}]
		Soit une partie \( A\) non vide dans \( \eN\). Si \( 0\in A\), nous avons fini.
		\begin{subproof}
			\spitem[L'ensemble \( B\)]
			Nous supposons donc que \( A\) ne contient pas zéro et nous définissons
			\begin{equation}
				B=\{ n\in \eN\setminus A\tq n\leq a, \forall a\in A \}.
			\end{equation}
			\spitem[Un élément particulier dans \( B\)]
			L'ensemble \( B\) vérifie :
			\begin{itemize}
				\item \( 0\in B\)
				\item \( B\neq \eN\) parce que \( A\) est non vide.
			\end{itemize}
			La contraposée de la condition \ref{ITEMooXPYEooFajywh} de la définition \ref{DEFooBJBOooWlblAx} d'un triplet naturel implique que \( s(B)\nsubset B\). Autrement dit, il existe \( b\in B\) tel que \( s(b)\notin B\).

			\spitem[Deux fonctions sur \( A\)]
			Puisque \( b\in B\), nous avons une application \( c\colon A\to \eN\) telle que \( b+c(a)=a\). Nous avons \( c(a)\neq 0\) parce que \( c(a)=0\) signifierait \( b=a\), ce qui est impossible parce que \( a\in A\) et \( b\in B\).

			Comme pour tout \( a\in A\), l'élément \( c(a)\) est non nul, il existe une fonction \( d\colon A\to \eN\) telle que \( c(a)=d(a)+1\).
			\spitem[\( s(b)\in A\)]
			Supposons que \( s(b)\notin A\). Alors il existe \( a\in A\) tel que \( s(b)\leq a\) est faux. Puisque l'ordre est total (proposition \ref{PROPooGCCRooFBYrlo}), nous avons
			\begin{equation}        \label{EQooYQSFooPSPJMt}
				a\leq s(b).
			\end{equation}
			Comme \( b\in B\) nous avons aussi
			\begin{equation}        \label{EQooIPAWooDBSJEa}
				b\leq a.
			\end{equation}
			Et enfin nous avons
			\begin{equation}        \label{EQooWFHRooGDSBFD}
				b\neq a
			\end{equation}
			parce que \( a\in A\) et \( b\in B\).

			Les conditions \eqref{EQooIPAWooDBSJEa} et \eqref{EQooWFHRooGDSBFD} se résument en \( b<a\). Le lemme \ref{LEMooPVRQooXPMKTt}\ref{ITEMooWCOIooMWrCag} nous indique alors que \( s(b)\leq a\). Cela mis à côté de \eqref{EQooYQSFooPSPJMt} conclut que \( a=s(b)\), et donc que \( s(b)\) est dans \( A\). Contradiction. Nous en concluons que \( s(b)\in A\).
			\spitem[\( s(b)\) est un minimum de \( A\)]
			En utilisant la commutativité et l'associativité de la somme nous avons, pour tout \( a\in A\) :
			\begin{equation}
				a=b+c(a)=b+\big( d(a)+1 \big)=(b+1)+d(a)=s(b)+d(a).
			\end{equation}
			Donc \( s(b)\leq a\) pour tout \( a\in A\). Mais comme \( s(b)\in A\), l'élément \( s(b)\) est bien un minimum de \( A\).
			\spitem[Unicité]
			Si \( a\) et \( a'\) sont des minimums de \( A\), alors \( a\leq a'\) et \( a'\leq a\). Nous en déduisons que \( a=a'\).
		\end{subproof}
		\spitem[Pour \ref{ITEMooSRGOooNYJJHY}]
		Si \( M\in \eN\) majore tous les éléments de \( \eN\), alors en particulier \( M\geq s(M)\). Mais le lemme \ref{LEMooWMYPooLTMyWb} nous indique que \( M\leq s(M)\). Nous avons donc \( s(M)=M\), c'est-à-dire \( M=M+1\). En utilisant la régularité de \( M\)\footnote{Dit plus simplement : en simplifiant par \( M\).}, nous trouvons \( 0=1\), ce qui est impossible par le lemme \ref{LEMooCOMSooEWrumL}.
		\spitem[Pour \ref{ITEMooKIHZooDRTCdx}]
		Soit \( A\) une partie non vide et majorée de \( \eN\).
		\begin{subproof}
			\spitem[L'ensemble \( B\)]
			Nous posons
			\begin{equation}
				B=\{ n\in \eN\setminus A\tq a\leq n, \forall a \in A \}.
			\end{equation}
			\spitem[\( B\) est non vide]
			Soit un majorant \( M\) de \( A\) : pour tout \( a\in A\) nous avons \( a\leq M\). Nous avons \( s(M)\notin A\), parce que si \( s(M)\) était dans \( A\), ce serait un élément de \( A\) strictement plus grand que tout \( a\in A\). Donc \( B\) est non vide parce qu'il contient \( s(M)\).
		\end{subproof}
		\spitem[Minimum]
		Puisque \( B\) est non vide, il possède un plus petit élément que nous notons \( b\). Nous savons que \( b\neq 0\) parce que sinon \( A\) serait vide. Il existe donc \( c\in \eN\) tel que \( b=s(c)\).

		\spitem[\( a\leq c\) pour tout \( a\in A\)]
		Comme \( s(c)\in B\) nous avons \( a< s(c)\) pour tout \( a\in A\). Donc, par le lemme \ref{LEMooPVRQooXPMKTt}\ref{ITEMooWCOIooMWrCag} nous avons \( s(a)\leq s(c)\), c'est-à-dire \( a+1\leq c+1\). Par régularité nous avons \( a\leq c\).

		Nous avons prouvé que \( a\leq c\) pour tout \( a\in A\).
		\spitem[\( c\in A\)]
		Si \( c\) n'est pas dans \( A\), alors il est dans \( B\) et il contredit la minimalité de \( b\). Donc \( c\) est dans \( A\).
		\spitem[Conclusion]
		L'élément \( c\) est dans \( A\) tout en étant plus petit que tout élément de \( A\).
		\spitem[Unicité]
		Si \( x\) est un élément minium de \( A\), alors nous avons \( x\leq c\) parce que \( x\) est minimum et \( c\leq x\) parce que \( c\) est minimum, et donc \( x=c\).
	\end{subproof}
\end{proof}

\begin{proposition}[\cite{MonCerveau}]	\label{PROPooTQDJooFQGFpC}
	Soient \( s,k\in \eN\).
	\begin{enumerate}
		\item		\label{ITEMooPQETooJYjhUj}
		      Si \( s>k\) alors \( s-1\geq k\).
		\item		\label{ITEMooZLXQooPRLNoN}
		      Si \( s\geq k\) alors \( s+1>k\).
	\end{enumerate}
\end{proposition}

\begin{proof}
	En deux parties.
	\begin{subproof}
		\spitem[Pour \ref{ITEMooPQETooJYjhUj}]
		%-----------------------------------------------------------
		Nous supposons que \( s>k\). Il existe \( x\in \eN\) tel que \( s+x=k\) et \( x\neq 0\) parce que \( s\neq k\). En utilisant la proposition \ref{PROPooAKLVooBkRIAT} nous avons \( k+(x-1)=s-1\), c'est-à-dire \( k\geq s-1\).
		\spitem[Pour \ref{ITEMooZLXQooPRLNoN}]
		%-----------------------------------------------------------
		Nous savons que \( s+x=k\) pour un certain \( x\in \eN\). Nous avons alors
		\begin{equation}
			(s+1)+x=k+1>k.
		\end{equation}
	\end{subproof}
\end{proof}



\begin{probleme}		\label{PROBooUCQFooTugfDl}
	Je ne suis pas certain qu'il y ait moyen de prouver le lemme \ref{LEMooPCRFooXRGrUr} sans d'abord définir les naturels. Une raison est que la véracité de ce lemme est à peu près équivalente au théorème \ref{THOooEJPYooZFVnez} qui donne l'existence de \( \eN\).

	J'ai essayé un certain nombre de choses avec le lemme de Zorn sans succès.

	Si vous avez une idée, n'hésitez pas à me la partager, ou à l'écrire ici :
	\begin{center}
		\url{https://math.stackexchange.com/questions/5019157/every-non-empty-dedekind-finite-totally-ordered-set-has-a-maximum-using-zorns}
	\end{center}
	%TODOooODPBooJZFYau. Voir si il y a une réponse. C'est déjà dans ma liste.

	Dans la réponse de Asaf Karagila, il y a seulement le fait que \(f \colon D\to D  \) que je ne comprends pas. C'est d'ailleurs ma dernière question dans les commentaires.
\end{probleme}




\begin{lemma}[\cite{MonCerveau}]       \label{LEMooYMRJooYIAhBb}
	Quelques affirmations sur l'ordre dans \( \eN\).
	\begin{enumerate}
		\item       \label{ITEMooTLOIooTWNtod}
		      Il n'existe pas de \( n\in \eN\) tel que \( n<0\).
		\item       \label{ITEMooPJKQooGfLCUM}
		      Si \( a,b\in \eN\) vérifient \( a>b\), alors il n'existe pas de \( x\) dans \( \eN\) tel que \( a+x=b\).
	\end{enumerate}
\end{lemma}

\begin{proof}
	En deux parties.
	\begin{subproof}
		\spitem[Pour \ref{ITEMooTLOIooTWNtod}]
		Nous savons par la proposition \ref{PROPooMZOWooHmsXzI}\ref{ITEMooJLAHooDKukfH} que \( 0\) est l'unique minimum de \( \eN\). Nous avons donc forcément \( 0\leq n\). Si \( n\) vérifie de plus \( n\leq 0\) alors nous avons \( n=0\) par symétrie de la relation d'ordre \( \leq\). Il n'est donc pas possible d'avoir \( n\neq 0\).
		\spitem[Pour \ref{ITEMooPJKQooGfLCUM}]
		Si \( b\leq a\) il existe \( c\in \eN\) tel que \( b+c=a\). Et comme \( a\neq b\), \( c\) n'est pas nul et il existe \( y\in \eN\) tel que \( c=s(y)\). Bref, nous avons
		\begin{equation}
			b+s(y)=a.
		\end{equation}
		Si de plus il existe \( x\in \eN\) tel que \( a+x=b\) nous aurions
		\begin{equation}
			a+x+s(y)=a.
		\end{equation}
		Comme \( a\) est régulier pour l'addition\footnote{Proposition \ref{PROPooTLTSooGNMTmV}\ref{ITEMooNUTHooJWWzGv}.}, nous avons
		\begin{equation}
			x+s(y)=0,
		\end{equation}
		ce qui signifie, par le lemme \ref{LEMooQBHFooCuCusQ} que \( x=s(y)=0\). Puisque \( s\) prend ses valeurs dans \( \eN\setminus\{ 0 \}\), cela est impossible.
	\end{subproof}
\end{proof}

\begin{definition}      \label{DEFooKBUFooLvMHrf}
	Soient \( a,b\in \eN\) tels que \( a\leq b\). Nous notons par \( \{ a,\ldots, b \}\) l'ensemble
	\begin{equation}
		\{ x\in \eN\tq a\leq x\leq b \}.
	\end{equation}
\end{definition}

\begin{proposition}     \label{PROPooFYMJooWihvhk}
	Toute application \( \eN\to \eN\) strictement croissante est injective.
\end{proposition}

\begin{proof}
	Soit une application strictement croissante \( f\colon \eN\to \eN\). Soient \( a,b\in \eN\) tels que \( f(a)=f(b)\). Puisque l'ordre est total\footnote{Proposition \ref{PROPooGCCRooFBYrlo}.}, nous supposons que \( a\leq b\). Si \( a=b\) nous avons terminé. Nous supposons donc que \( a\neq b\), c'est-à-dire que \( a<b\). Par stricte croissance nous avons alors \( f(a)<f(b)\) qui signifie \( f(a)\leq f(b)\) et \( f(a)\neq f(b)\). Contradiction. Il n'existe donc pas de \( a\neq b\) tels que \( f(a)=f(b)\). L'application \( f\) est donc injective.
\end{proof}

\begin{lemma}[\cite{MonCerveau}]        \label{LEMooFKLPooPrmeUU}
	Si \( S\) n'est pas majoré dans \( \eN\), alors il existe une bijection \( \eN\to S\).
\end{lemma}

\begin{proof}
	Nous considérons l'application suivante :
	\begin{equation}
		\begin{aligned}
			g\colon S & \to S                            \\
			n         & \mapsto \min\{ x\in S\tq x>n \}.
		\end{aligned}
	\end{equation}
	Cette application est bien définie parce que tout partie non vide de \( \eN\) a un plus petit élément\footnote{Proposition \ref{PROPooMZOWooHmsXzI}\ref{ITEMooYAJIooEFmOpB}.}. Maintenant nous définissons \( f\colon \eN\to S\) par
	\begin{subequations}
		\begin{numcases}{}
			f(0)=\min(S)\\
			f(n+1)=g\big( f(n) \big).
		\end{numcases}
	\end{subequations}
	C'est le théorème \ref{THOooEJPYooZFVnez} qui nous permet de le faire. Nous montrons que \( f\) est bijective.
	\begin{subproof}
		\spitem[Injective]
		Nous avons
		\begin{equation}
			f(n+1)\in\{ x\in S\tq x>f(n) \}.
		\end{equation}
		Donc \( f\) est strictement croissante. Elle est donc injective.
		\spitem[Surjective]
		Soit \( a\in S\). Nous allons voir que \( a\) est dans l'image de \( f\). Pour cela nous posons
		\begin{equation}
			A=\{ x\in \eN\tq f(x)<a \}.
		\end{equation}
		Cet ensemble est majoré par \( a\). En effet si \( x\in A\) nous avons \( x\leq f(x)<a\). La partie \( A\) de \( \eN\) possède un maximum. Nous notons \( M=\max(A)\). Ce \( M\) a deux propriétés intéressantes.
		\begin{subproof}
			\spitem[D'abord]
			Puisque \( M\in A\), nous avons \( f(M)<a\). Une autre façon de dire cela est de dire que
			\begin{equation}
				a\in \{ x\in S\tq x>f(M) \}.
			\end{equation}
			Or \( f(M+1)=\min\{ x\in S\tq x>f(M) \}\). Donc \( f(M+1)\leq a\).
			\spitem[Ensuite]
			Puisque \( M\) est le maximum de \( A\), \( M+1\) majore \( A\), c'est-à-dire que \( f(M+1)\geq a\).
			\spitem[Les deux ensemble]
			Nous avons prouvé que \( f(M+1)\leq a\) et \( f(M+1)\geq a\). Nous en déduisons, par le lemme \ref{LEMooJRZKooOMhOkH}, que \( f(M+1)=a\).
		\end{subproof}
	\end{subproof}
\end{proof}

\begin{normaltext}      \label{NORMooQXASooMXqhjI}
	Durant la preuve du lemme \ref{LEMooFKLPooPrmeUU}, nous n'avons pas été loin de prouver que
	\begin{equation}
		\big( \min(S),S,g \big)
	\end{equation}
	est un triplet naturel.

	Toute partie non bornée de \( \eN\) donne lieu à un triplet naturel.
\end{normaltext}



\begin{lemma}[\cite{MonCerveau}]	\label{LEMooLDKQooSTiVlW}
	Si \( A\subset \eN\), alors
	\begin{equation}
		\min\big( A\setminus\{ \min(A) \} \big)>\min(A).
	\end{equation}
\end{lemma}

\begin{proof}
	Soient \( m=\min(A)\) et \( m'=\min\big( A\setminus\{ m \} \big)\). Il y a trois possibilités : \( m>m'\), \( m=m'\) et \( m<m'\). Supposons que \( m'\leq m\). Étant donné que \( m'\in A\), nous avons \( m\leq m'\). Et donc \( m=m'\), ce qui est impossibe parce que \( m'\in A\setminus\{ m \}\).
\end{proof}


%---------------------------------------------------------------------------------------------------------------------------
\subsection{Multiplication dans les naturels}
\label{SUBooMultiplicationNaturels}
%---------------------------------------------------------------------------------------------------------------------------

\begin{definition}[\cite{MonCerveau}]      \label{DEFooAZAYooVjNzmy}
	Soit un ensemble muni d'une loi de composition interne \( (A,+)\). Soit \( n\in \eN\) et \( a\in A\). Nous définissons \( n\times A\) par
	\begin{subequations}
		\begin{numcases}{}
			0\times a=0\\
			(n+1)\times a= n\times a+a.
		\end{numcases}
	\end{subequations}
\end{definition}

\begin{definition}      \label{DEFooLCWLooYrToFv}
	Un ensemble totalement ordonné muni d'une loi de composition interne \( (A,+, \leq)\) est \defe{archimédien}{ensemble!archimédien}\index{archimédien} si pour tout \( x,y\in A\) avec \( x>0\), il existe \( n\in \eN\) tel que \( n\times x\geq y\) (voir la définition \ref{DEFooAZAYooVjNzmy}).
\end{definition}



\begin{propositionDef}      \label{PROPooBBQPooRgPOjf}
	Il existe une unique fonction \( f\colon \eN\times \eN\to \eN\) telle que
	\begin{enumerate}
		\item       \label{ITEMooNTUUooDAUVsV}
		      \( f(a,0)=0\) pour tout \( a\in \eN\)
		\item       \label{ITEMooPPZZooQQabwn}
		      \( f\big( a,s(b) \big)=f(a,b)+a\) pour tout \( a,b\in \eN\).
	\end{enumerate}
	Cette fonction est la \defe{multiplication}{multiplication} et nous notons \( f(a,b)=a\times b\), voire \( ab\) quand il n'y a pas d'ambigüité. Le nombre \( a\times b\) est nommé le \defe{produit}{produit} de \( a\) par \( b\).
\end{propositionDef}

\begin{proof}
	En deux parties.
	\begin{subproof}
		\spitem[Fonctions définies par récurrence]
		Soit \( a\in \eN\). Par le théorème \ref{THOooEJPYooZFVnez}, il existe une unique application \( f_a\colon \eN\to \eN\) telle que
		\begin{subequations}        \label{EQSooWVCTooNTVjKU}
			\begin{numcases}{}
				f_a(0)=0\\
				f_a\big( s(b) \big)=f_a(b)+a
			\end{numcases}
		\end{subequations}
		\spitem[Existence]
		Nous considérons, pour chaque \( a\in \eN\) la fonction \( f_a\) définie par les conditions \eqref{EQSooWVCTooNTVjKU}. En posant \( f(a,b)=f_a(b)\), nous avons une application qui vérifie toutes les conditions.
		\spitem[Unicité]
		Soient des applications \( f\) et \( g\) vérifiant les propriétés demandées. Soit \( a\in \eN\). Nous pouvons définir \( f_a\colon \eN\to \eN\) et \( g_a\colon \eN\to \eN\) par \( f_a(n)=f(a,b)\) et \( g_a(b)=g(a,b)\).

		Les applications \( f_a\) et \( g_a\) vérifient toutes deux les conditions \eqref{EQSooWVCTooNTVjKU}, et sont donc égales : \( f_a=g_a\) pour tout \( a\). Donc \( f=g\).
	\end{subproof}
\end{proof}

\begin{normaltext}	\label{NORMooPrioriteOperations}
	Nous supposons que \randomGender{le lecteur}{la lectrice} connait déjà la priorité des opérations. \randomGender{Il}{Elle} saura donc interpréter des expressions comme \( a\times b+c\) comme voulant dire \( (a\times b)+c\) sans que nous ayons à ajouter des parenthèses.
\end{normaltext}

\begin{proposition}[\cite{RWWJooJdjxEK,MonCerveau}]     \label{PROPooGHDOooFYRmon}
	La multiplication dans \( \eN \) possède les propriétés suivantes.
	\begin{enumerate}
		\item       \label{ITEMooHFWRooDCEpjj}
		      \( n\times 1=n\) pour tout \( n\in \eN\).
		\item       \label{ITEMooRSYMooSUrRsl}
		      \( 1\times n=n\) pour tout \( n\in \eN\).
		\item       \label{ITEMooWJPOooRUYjwQ}
		      La multiplication est commutative.
		\item       \label{ITEMooNBYKooXnGRrf}
		      \( 0\times n=0\) pour tout \( n\in \eN\)
		\item      \label{ITEMooLJQBooVpUxUv}
		      L'élément \( 1\) est neutre pour la multiplication.
		\item       \label{ITEMooDYLIooETIBEL}
		      La multiplication est distributive par rapport à l'addition.
		\item       \label{ITEMooQBFSooWGDQYX}
		      La multiplication est associative.
	\end{enumerate}
\end{proposition}

\begin{proof}
	Point par point.
	\begin{subproof}
		\spitem[Pour \ref{ITEMooHFWRooDCEpjj}]
		Nous avons \( n\times 1=n\times s(0)=(n\times 0)+n=0+n=n\). Donc \( n\times 1=n\).
		\spitem[Pour \ref{ITEMooRSYMooSUrRsl}]
		Nous le faisons par récurrence. Par définition c'est vrai pour \( n=0\). En ce qui concerne la récurrence, nous supposons que \( 1\times n=n\), et nous prouvons que \( 1\times s(n)=s(n)\) :
		\begin{equation}
			1\times s(n)=(1\times n)+1=n+1=s(n).
		\end{equation}
		\spitem[Pour \ref{ITEMooWJPOooRUYjwQ}]
		Soit \( a\in \eN\). Nous prouvons par récurrence sur \( b\in \eN\) que \( a\times b=b\times a\). Pour \( b=0\) c'est bon. Pour la récurrence, nous supposons que \( a\times b=b\times a\) et nous prouvons que \( a\times s(b)=s(b)\times a\) :
		\begin{subequations}
			\begin{align}
				a\times s(b) & =a\times b+a                                                   \\
				             & =b\times a+a             & \text{récurrence}                   \\
				             & =(b\times a)+(1\times a) & \text{par \ref{ITEMooRSYMooSUrRsl}} \\
				             & =(b+1)\times a           & \text{distributivité}               \\
				             & =s(b)\times a.
			\end{align}
		\end{subequations}
		\spitem[Pour \ref{ITEMooNBYKooXnGRrf}]
		C'est vrai pour \( n=0\) par la définition \ref{PROPooBBQPooRgPOjf}\ref{ITEMooNTUUooDAUVsV}. En ce qui concerne \( s(n)\), nous avons
		\begin{equation}
			0\times s(n)=(0\times n)+0=0.
		\end{equation}
		\spitem[Pour \ref{ITEMooLJQBooVpUxUv}]
		C'est la combinaison de \ref{ITEMooHFWRooDCEpjj} et \ref{ITEMooRSYMooSUrRsl}.
		\spitem[Pour \ref{ITEMooDYLIooETIBEL}]
		Soient \( a,b\in \eN\). Nous prouvons par récurrence sur \( c\in \eN\) que\footnote{Nous n'écrivons pas toutes les parenthèses parce que les règles de priorité des opérations sont supposées connues. J'invite cependant \randomGender{le lecteur}{la lectrice} à remarquer qu'une formalisation de ces règles n'est probablement pas facile. Pour que tout soit rigoureux, il faudrait un algorithme qui parcourt une suite de caractères et l'interprète en ajoutant correctement les parenthèses.}
		\begin{equation}
			(a+b)\times c=a\times c+b\times c.
		\end{equation}
		Pour \( c=0\), nous avons \( (a+b)\times 0=0\) ainsi que \( a\times 0=b\times 0=0\) en vertu des points précédents sur la multiplication par zéro. Pour la récurrence nous utilisons associativité et commutativité de la somme :
		\begin{subequations}
			\begin{align}
				(a+b)\times s(c) & =(a+b)\times c+(a+b)                                \\
				                 & =a\times c+b\times c+a+b                            \\
				                 & =(a\times c +a)+(b\times c+b)                       \\
				                 & =\big( a\times s(c) \big)+\big( b\times s(c) \big).
			\end{align}
		\end{subequations}
		On devrait aussi prouver que \( c(a+b) = ca + cb \): c'est peu ou prou la même chose.
		\spitem[Pour \ref{ITEMooQBFSooWGDQYX}]
		Soient \( a,b\in \eN\). Nous démontrons par récurrence sur \( c\in \eN\) que \( (a\times b)\times c=a\times (b\times c)\). Pour \( c=0\) l'égalité est triviale. Nous supposons que l'égalité est correcte pour \( c\), et nous la prouvons pour \( s(c)\) :
		\begin{subequations}
			\begin{align}
				(a\times b)\times s(c) & =\big( (a\times b)\times c \big)+a\times b                                                                      \\
				                       & =\big( a\times (b\times c) \big)+a\times b & \text{récurrence}                                                  \\
				                       & =\big( (b\times c)\times a \big)+b\times a & \text{commutativité}                                               \\
				                       & =(b\times c+b)\times a                     & \text{distributivité}                                              \\
				                       & =\big( b\times s(c) \big)\times a          & \text{définition \ref{PROPooBBQPooRgPOjf}\ref{ITEMooPPZZooQQabwn}} \\
				                       & =a\times \big( b\times s(c) \big)          & \text{commutativité}.
			\end{align}
		\end{subequations}
	\end{subproof}
\end{proof}

\begin{lemma}[\cite{MonCerveau}]	\label{LEMooCompatMultiplicationOrdre}
	La multiplication est compatible avec l'ordre :
	\begin{equation}
		a\leq b\Rightarrow a\times n\leq b\times n
	\end{equation}
	pour tout \( n\in \eN\).
\end{lemma}

\begin{proof}
	Par la définition \ref{DEFooAXZSooTEMjlV} de l'ordre, si \( a\leq b\), il existe \( c\in \eN\) tel que \( b=a+c\). En utilisant la distributivité\footnote{Proposition \ref{PROPooGHDOooFYRmon}\ref{ITEMooDYLIooETIBEL}.}, nous avons
	\begin{equation}
		b\times n=(a+c)\times n=a\times n+c\times n.
	\end{equation}
	Nous en déduisons que \( a\times n\leq b\times n\) parce que \( c\times n \in \eN\).
\end{proof}

\begin{lemma}       \label{LEMooEHYEooLDudfn}
	Si \( a\times b=0\), alors \( a=0\) ou \( b=0\) (ou les deux).
\end{lemma}

\begin{proof}
	Supposons que \( a\neq 0\). Alors il existe \( c\in \eN\) tel que \( a=s(c)\). Nous avons
	\begin{equation}
		0=a\times b=s(c)\times b=c\times b+b.
	\end{equation}
	Le lemme \ref{LEMooQBHFooCuCusQ} nous dit alors que \( c\times b=b=0\).
\end{proof}

\begin{lemma}[\cite{MonCerveau}]        \label{LEMooGUXGooBcKJdS}
	Soit \(a, b, n\) trois entiers naturels. Si \( a< b\) et si \( n\neq 0\), alors
	\begin{equation}
		a\times n<b\times n.
	\end{equation}
\end{lemma}

\begin{proof}
	L'hypothèse \( a\leq b\) implique qu'il existe \( c\in \eN\) tel que \( a+c=b\). De plus \( c\neq 0\) parce que \( a\neq b\). En utilisant la distributivité\footnote{Proposition \ref{PROPooGHDOooFYRmon}\ref{ITEMooDYLIooETIBEL}.}, nous avons
	\begin{equation}
		b\times n=(a+c)\times n=(a\times n)+(c\times n).
	\end{equation}
	Cela prouve que \( a\times n\leq b\times n\). Et comme \( c\) et \( n\) ne sont pas nuls, nous avons même\footnote{Lemme \ref{LEMooEHYEooLDudfn}.} \( c\times n\neq 0\) et donc \( a\times n<b\times n\).
\end{proof}

\begin{remark}	\label{REMooultiplicationNaturelsOrdreStrictZ}
	Une version du lemme \ref{LEMooGUXGooBcKJdS} dans \( \eZ\) sera le lemme \ref{LEMooSVDDooWsyxNP}.
\end{remark}

\begin{lemma}[\cite{MonCerveau}]        \label{LEMooSFUKooBNAple}
	Soient \( a\neq 0\) et \( b>1\) dans \( \eN\). Alors
	\begin{equation}
		ab>a.
	\end{equation}
\end{lemma}

\begin{proof}
	Il s'agit d'une application du lemme \ref{LEMooGUXGooBcKJdS} en partant de l'inégalité \( 1<b\) et en la «multipliant» par \( a\).
\end{proof}

\begin{proposition}[\cite{MonCerveau}]		\label{PROPooNaturelsReguliersMultiplication}
	Tous les naturels non nuls sont réguliers par rapport à la multiplication. Autrement dit, si \( a\neq 0\), alors nous avons
	\begin{equation}
		a\times x=a\times y\Rightarrow x=y.
	\end{equation}
\end{proposition}

\begin{proof}
	Soit \( a\neq 0\) dans \( \eN\). Nous supposons que \( a\times x=a\times y\). Puisque l'ordre sur \( \eN\) est total (proposition \ref{PROPooGCCRooFBYrlo}), nous pouvons supposer que \( y\geq x\); sinon il suffit de permuter les rôles de \( x\) et \( y\) dans tout ce qui suit.

	Il existe \( d\in \eN\) tel que \( y=x+d\). En utilisant l'hypothèse \( a\times y=a\times x\) et la distributivité\footnote{Proposition \ref{PROPooGHDOooFYRmon}\ref{ITEMooDYLIooETIBEL}.},
	\begin{equation}
		a\times x=a\times y=a\times (x+d)=(a\times x)+(a\times d).
	\end{equation}
	Puisque \( (a\times x)\) est régulier pour la somme\footnote{Proposition \ref{PROPooTLTSooGNMTmV}\ref{ITEMooNUTHooJWWzGv}.} nous en déduisons que
	\begin{equation}
		0=a\times d.
	\end{equation}
	Le lemme \ref{LEMooEHYEooLDudfn} dit alors que \( a=0\) ou que \( d=0\). Étant donné que \( a\neq 0\) par hypothèse, nous déduisons que \( d=0\), c'est-à-dire que \( x=y\).
\end{proof}

%---------------------------------------------------------------------------------------------------------------------------
\subsection{Presque unicité des triplets naturels}
\label{SUBooUniciteTripletsNaturels}
%---------------------------------------------------------------------------------------------------------------------------

\begin{normaltext}	\label{NORooIntroUniciteTripletsNaturels}
	Il existe de nombreux triplets naturels; l'existence d'un triplet naturel est un théorème de la théorie des ensembles que nous avons accepté. Nous avons déjà à peu près montré que toute partie non bornée de \( \eN\) donne lieu à un nouveau triplet naturel. Voir \ref{NORMooQXASooMXqhjI}.

	Nous voyons maintenant que tous les triplets naturels sont équivalents au moins pour l'ordre.
\end{normaltext}

\begin{theorem}[\cite{MonCerveau}]     \label{THOooFUXMooJuigHK}
	Soient des triplets naturels \( (\mN_1,o_1,s_1)\) et \( (\mN_2,o_2,s_2)\). Alors
	\begin{enumerate}
		\item
		      il existe une unique application \( f\colon \mN_1\to \mN_2\) telle que
		      \begin{enumerate}
			      \item
			            \( f(o_1)=o_2\)
			      \item
			            \( f\circ s_1=s_2\circ f\).
		      \end{enumerate}
		\item
		      Une telle application est une bijection croissante.
	\end{enumerate}
\end{theorem}

\begin{proof}
	En plusieurs points.
	\begin{subproof}
		\spitem[Existence]
		Nous voyons \( (\mN_1, o_1, s_1)\) comme un triplet naturel, et \( \mN_2\) comme un simple ensemble. Nous pouvons appliquer le théorème \ref{THOooEJPYooZFVnez} à \( (\mN_1, o_1, s_1)\). L'élément \( o_1\) va jouer le rôle de \( 0\) alors que \( o_2\) va jouer le rôle de \( b\). L'application \( g\) est \( s_2\). Bref, il existe une unique application \( f\colon \mN_1\to \mN_2\) telle que
		\begin{enumerate}
			\item
			      \( f(o_1)=o_2\)
			\item
			      \( f\big( s_1(n) \big)=s_2\big( f(n) \big)\)
		\end{enumerate}
		pour tout \( n\in \mN_1\).
		\spitem[Unicité]
		Le théorème \ref{THOooEJPYooZFVnez} donne déjà l'unicité. Nous la faisons quand même, juste pour vous faire plaisir. Soit \( g\), une autre application vérifiant les mêmes conditions. Pour faire la récurrence de façon très explicite, nous posons
		\begin{equation}
			\begin{aligned}
				P\colon \mN_1 & \to \{ 0,1 \}                    \\
				x             & \mapsto \begin{cases}
					                        1 & \text{si } g(x)=f(x) \\
					                        0 & \text{sinon. }
				                        \end{cases}
			\end{aligned}
		\end{equation}
		Notre but est de prouver que \( P(x)=1\) pour tout \( x\in \mN_1\), en utilisant la récurrence telle que décrite dans la proposition \ref{PROPooXTRCooKwrWkq}.

		Nous avons \( f(o_1)=o_2=g(o_1)\). Donc \( P(o_1)=1\). Nous supposons que, pour un certain \( a\in \mN_1\), nous ayons \( P(a)=1\), et nous prouvons que \( P\big( s_1(a) \big)=1\).

		Nous avons \( g(a)=f(a)\), et nous prenons \( s_2\) des deux côtés, nous avons succéssivement
		\begin{subequations}
			\begin{align}
				(s_2\circ g)(a)     & =(s_2\circ f)(a)      \\
				(g\circ s_1)(a)     & =(f\circ s_1)(a)      \\
				g\big( s_1(a) \big) & =f\big( s_1(a) \big).
			\end{align}
		\end{subequations}
		La dernière égalité signifie que \( P\big( s_1(a) \big)=1\). La proposition \ref{PROPooXTRCooKwrWkq} implique que \( P(x)=1\) pour tout \( x\in \mN_1\).
		\spitem[Bijection, définir l'inverse]
		Nous allons trouver un inverse et le lemme \ref{LEMooWBYSooFqyqQP} nous dit que c'est suffisant. La partie «existence», en inversant les rôles de \( \mN_1\) et \( \mN_2\) nous donne une application \( g\colon \mN_2\to \mN_1\) telle que
		\begin{enumerate}
			\item
			      \( g(o_2)=o_1\)
			\item
			      \( g\circ s_2=s_1\circ g\).
		\end{enumerate}
		Nous allons prouver que \( g\) est un inverse de \( f\).
		\spitem[\( f\circ g=\id\)]
		Nous posons \( A=\{ x\in \mN_2\tq (f\circ g)(x)=x \}\). Nous avons
		\begin{equation}
			f\big( g(o_2) \big)=f(o_1)=o_2,
		\end{equation}
		et donc \( o_2\in A\).

		Supposons que \( x\in A\). Alors
		\begin{subequations}
			\begin{align}
				(f\circ g)\big( s_2(x) \big) & =(f\circ \underbrace{g\circ s_2}_{s_1\circ g})(x)  \\
				                             & =(\underbrace{f\circ s_1}_{=s_2\circ f}\circ g)(x) \\
				                             & =(s_2\circ f\circ g)(x)                            \\
				                             & =s_2\big( (f\circ g)(x) \big)                      \\
				                             & =s_2(x)
			\end{align}
		\end{subequations}
		Donc \( s_2(x)\in A\). Nous en déduisons que \( A=\mN_2\) par le point \ref{ITEMooXPYEooFajywh} de la définition \ref{DEFooBJBOooWlblAx} d'un triplet naturel.
		\spitem[\( g\circ f=\id\)]
		J'imagine que c'est la même chose que dans l'autre sens (ci-dessus)\quext{Je n'ai pas fait les calculs; écrivez-moi si ça pose un problème.}.
	\end{subproof}
\end{proof}

\begin{proposition}     \label{PROPooCCVNooYUYcqG}
	L'ensemble structuré \( (\eN,+,\times, \leq)\) est archimédien\footnote{Définition \ref{DEFooLCWLooYrToFv}.}. En d'autres termes, pour tout \( a,b\in \eN\setminus\{ 0 \}\), il existe \( n\in \eN\) tel que
	\begin{equation}
		b<n\times a.
	\end{equation}
	\index{archimédien!\( \eN\)}
\end{proposition}

\begin{proof}
	Soient \( a,b\in \eN\) avec \( a\neq 0\).

	Si \( a>b\), nous avons le résultat avec \( n=1\).

	Si \( a=b\), en prenant \( n=s(1)\) nous avons le résultat. En effet \( s(1)\times a=a+a\). Puisque \( a\neq 0\), nous avons \( a+a\geq a\) et \( a+a\neq a\), donc \( s(1)\times a>a\).

	La vraie vie est avec \( a<b\). Nous posons
	\begin{equation}
		X=\{ x\in \eN\tq 1\leq x\times a\leq b \}
	\end{equation}
	et
	\begin{equation}
		B=\{ x\times a\tq x\in X \}.
	\end{equation}
	L'ensemble \( X\) est non vide parce que \( 1\in X\). L'ensemble \( B\) est alors également non vide, et majoré par \( b\). La proposition \ref{PROPooMZOWooHmsXzI}\ref{ITEMooKIHZooDRTCdx} nous indique alors que \( B\) possède un plus grand élément que nous allons noter \( x_0\times a\) (\( x_0\in X\)).

	Nous posons \( n=s(x_0)\), et nous avons
	\begin{equation}
		x_0\times a< x_0\times a +a=s(x_0)\times a=n\times a.
	\end{equation}
	Nous en déduisons que \( n\times a\) n'est pas dans \( B\) parce que \( x_0\times a\) est le plus grand élément de \( B\). Donc \( x_0\) n'est pas dans \( X\); nous n'avons donc pas les inégalités
	\begin{equation}
		1\leq n\times a\leq b.
	\end{equation}
	Laquelle des deux inégalités est fausse ? Puisque \( n=s(x_0)\geq 1\) et que \( a\geq 1\), nous avons \( 1\leq n\times a\). Donc c'est la seconde inégalité qui est fausse. Nous avons donc \( n\times a>b\).
\end{proof}



%-------------------------------------------------------
\subsection{Cardinalité dans \( \eN\)}
%----------------------------------------------------

\begin{normaltext}	\label{NORMooEnsembleFiniEquivalenceDef}
	La proposition suivante est à peu près prise comme définition d'un ensemble fini dans \cite{ooVAYLooJxVYex} qui donne également une preuve de l'équivalence avec notre définition.
\end{normaltext}

\begin{proposition}     \label{PROPooBYKCooGDkfWy}
	L'ensemble \( \eN\) est infini\footnote{Définition \ref{DefEOZLooUMCzZR}.}.
\end{proposition}

\begin{proof}
	Nous considérons la partie propre\footnote{Le fait que ce soit une partie propre est dû au fait que \( 0\) n'est pas dedans d'une part parce que le lemme \ref{LEMooWMYPooLTMyWb} dit que \( 0\leq 1\), et d'autre part parce que le lemme \ref{LEMooQBHFooCuCusQ} donne \( 0\neq 1\).}
	\begin{equation}
		A=\{ n\in \eN\tq n\geq 1 \}.
	\end{equation}
	Ensuite nous posons
	\begin{equation}
		\begin{aligned}
			\sigma\colon \eN & \to A         \\
			x                & \mapsto s(x).
		\end{aligned}
	\end{equation}
	Le fait que \( \sigma\) prenne ses valeurs dans \( A\) est parce que \( s\) prend ses valeurs dans \( \eN\setminus\{ 0 \}=A\).
	\begin{subproof}
		\spitem[\( \sigma\) est injective]
		Parce que \( s\) l'est.
		\spitem[\( \sigma\) est surjective]
		C'est dans la définition \ref{DEFooBJBOooWlblAx}\ref{ITEMooQAKJooGKdJsM}
	\end{subproof}
	Nous avons donc une bijection entre \( \eN\) et un sous-ensemble strict.
\end{proof}

\begin{lemma}       \label{LEMooYIWDooIQFSad}
	Si \( a\leq b\) dans \( \eN\), alors la partie\footnote{Définition \ref{DEFooKBUFooLvMHrf}.} \( \{ a,\ldots, b \}\) est finie.
\end{lemma}

\begin{proof}
	En deux parties.
	\begin{subproof}
		\spitem[Pour \( a=0\)]
		%-----------------------------------------------------------
		Nous commençons par prouver que \( \{ 0,\ldots,b \}\) est fini. Par récurrence sur \( b\). Avec \( b=0\), la partie \( \{ 0 \}\) est finie parce que son seul sous-ensemble propre est \( \emptyset\) qui n'est pas en bijection avec \( \{ 0 \}\).

		Supposons que \( \{ 0,\ldots, b \}\) est fini. Alors \( \{ 0,\ldots, b \}\cup \{ b+1 \}\) est fini par le lemme \ref{PROPooVOKDooOStPzU}.

		\spitem[Cas général]
		%-----------------------------------------------------------
		Nous avons \( \{ a,\ldots,b \}\subset\{ 0,\ldots,b \}\). Donc \( \{ a,\ldots,b \}\) est une partie d'un ensemble fini. Elle est finie par \ref{LEMooTUIRooEXjfDY}.
	\end{subproof}
\end{proof}

\begin{lemma}       \label{LEMooJDGOooHdyJnu}
	Si \( p\neq q\), alors il n'existe pas de bijection entre \( \{ 0,\ldots, p \}\) et \( \{ 0,\ldots, q \}\).
\end{lemma}

\begin{proof}
	Supposons pour fixer les idées que \( p\leq q\). Dans ce cas \( \{ 0,\ldots, p \}\) est une partie stricte de \( \{ 0,\ldots, q \}\). Vu que \( \{ 0,\ldots, q \}\) est fini (lemme \ref{LEMooYIWDooIQFSad}), il n'y a pas de bijection avec ses parties strictes.
\end{proof}

\begin{proposition}[\cite{MonCerveau}]	\label{PROPooONYBooQuqAPR}
	Il n'existe pas de suites strictement décroissantes\footnote{Croissance, définition \ref{DEFooFTZAooLrtNEo}.} dans \( \eN\).
\end{proposition}

\begin{proof}
	Une suite dans \( \eN\) est une application \(f \colon \eN\to \eN  \). Nous supposons qu'elle est strictement décroissante. Étant donné que \( f(k)\leq f(0)\) pour tout \( k\), nous avons
	\begin{equation}
		f(\eN)\subset\{ x\in \eN\tq x\leq f(0) \}=\{ 0,\ldots,f(0) \}
	\end{equation}
	par la proposition \ref{PROPooHRCOooBkEvud}.

	D'autre part, l'ensemble \( \eN\) est infini par \ref{PROPooBYKCooGDkfWy}, et la partie \( \{ 0,\ldots,f(0) \}\) est finie par \ref{LEMooYIWDooIQFSad}. L'application \( f\) serait une injection d'un ensemble infini dans un ensemble fini, ce qui est impossible par la proposition \ref{PROPooWKSIooHcfYPN}.
\end{proof}


\begin{proposition}[\cite{RWWJooJdjxEK}]	\label{PROPooZLEXooGyjNRk}
	Toute partie non vide de \( \eN\) contient un unique élément minimal.
\end{proposition}

\begin{proof}
	Soit une partie non vide \( A\) de \( \eN\).
	\begin{subproof}
		\spitem[Unicité]
		%-----------------------------------------------------------
		Soient deux éléments minimaux \( m_1\) et \( m_2\). Si \( m_1\leq m_2\) alors \( m_1=m_2\) parce que \( m_2\) est minimal. Et si \( m_2\leq m_1\) alors \( m_1=m_2\) parce que \( m_1\) est minimal.
		\spitem[Existence]
		%-----------------------------------------------------------
		Supposons que \( A\) n'ait pas d'éléments minimaux. Vu qu'aucun élément de \( A\) est minimal, il existe une application \(f \colon A\to A  \) telle que \( f(a)<a\) pour tout \( a\). Nous considérons un élément \( x_0\in A\) et la suite définie par récurrence \( x_{n+1}=f(x_n)\) (via \ref{NORMooLNXMooWIblPf}).

		Cela donne une suite strictement décroissante dans \( \eN\), ce qui n'est pas possible par la proposition \ref{PROPooONYBooQuqAPR}.
	\end{subproof}
\end{proof}





%-------------------------------------------------------
\subsection{Division euclidienne}
%----------------------------------------------------


\begin{theorem}[division euclidienne \cite{RWWJooJdjxEK}]       \label{THOooKDJVooRIJRHP}
	Pour tout \( a\in \eN\), pour tout \( b\in \eN\setminus\{ 0 \}\), il existe un unique couple \( (q,r)\in \eN^2\) tel que \( a=bq+r\) avec \( 0\leq r<b\).
\end{theorem}

\begin{proof}
	Existence puis unicité.
	\begin{subproof}
		\spitem[Existence]

		Nous posons
		\begin{equation}
			A=\{ bx\tq x\in \eN,bx\leq a \}.
		\end{equation}
		L'ensemble \( A\) contient \( 0\) (avec \( x=0\)) et est majoré par \( a\). Donc il possède un plus grand élément que nous notons \( bq\). Puisque \( bq\in A\), nous avons \( bq\leq a\) et donc il existe \( r\in \eN\) tel que
		\begin{equation}        \label{EQooIUICooLenNBP}
			bq+r=a.
		\end{equation}
		Il reste à montrer que \( r<b\). Supposons que \( r\geq b\). Il existerait alors un \( x\) tel que \( b+x=r\). En mettant ça dans \eqref{EQooIUICooLenNBP},
		\begin{equation}
			bq+b+x=a,
		\end{equation}
		c'est-à-dire \( b(q+1)+x=a\), qui signifierait \( b(q+1)\leq a\), ce qui est faux parce que \( bq\) est le plus grand élément de \( A\).
		\spitem[Unicité]
		Supposons que nous ayons
		\begin{equation}
			a=bq+r=bq'+r'
		\end{equation}
		avec \( 0\leq r<b\) et \( 0\leq r'<b\). Il y a trois possibilités : \( q'<q\), \( q'=q\) et \( q'>q\).
		\begin{subproof}
			\spitem[Si \( q'<q\)]
			Alors il existe \( x\in \eN\) tel que \( q'+x=q\), et nous avons
			\begin{equation}
				b(q'+x)+r=bq'+bx+r,
			\end{equation}
			ce qui, après distribution et simplification, donne \( r'=bx+r\). Puisque nous avons \( x\geq 1\), il vient
			\begin{equation}
				r'=bx+r\geq b+r\geq b.
			\end{equation}
			Cela n'est pas possible parce que \( r'<b\). Le cas \( q'<q\) n'est pas possible.
			\spitem[Si \( q'=q\)]
			Nous avons alors immédiatement \( bq+r=bq+r'\) et donc \( r=r'\). Unicité.
			\spitem[Si \( q'>q\)]
			En posant \( q+x=q'\) nous trouvons la même impossibilité que dans le cas \( q'<q\).
		\end{subproof}
	\end{subproof}
\end{proof}

\begin{definition}	\label{DEFooDivisibleNombrePairImpair}
	Quelques définitions associées au théorème \ref{THOooKDJVooRIJRHP}. Soit \( a \) et \( b \) deux entiers naturels.
	\begin{enumerate}
		\item
		      Dans le couple \( (q,r)\in \eN^2\) tel que \( a=bq+r\) avec \( 0\leq r<b\), fourni par ce théorème, \( q \) est appelé \defe{quotient}{division euclidienne!quotient dans \( \eN \)} et \( r \) est appelé \defe{reste}{division euclidienne!reste dans \( \eN \)}.
		\item
		      Si \( r=0\) (et donc s'il existe un \( q \in \eN\) tel que \( a = bq \)), nous disons que \( a\) est \defe{divisible}{nombre entier!divisible} par \( b\).
		\item
		      Si \( n\) est divisible par \( 2\), alors nous disons qu'il est \defe{pair}{nombre entier!pair}.
		\item\label{ITEMooDLHHooTmWSnN}
		      Si \( n\) n'est pas divisible par \( 2\), alors nous disons qu'il est \defe{impair}{nombre entier!impair}.
	\end{enumerate}
\end{definition}


%+++++++++++++++++++++++++++++++++++++++++++++++++++++++
\section{Ensembles finis}
%+++++++++++++++++++++++++++++++++++++++++++++++++++++++

Voir \ref{NORMooLNLHooLYKwDu} pour comprendre le type de résultats que nous groupons ici.

\begin{lemma}[\cite{MonCerveau}]	\label{LEMooKNKEooYQeDdo}
	Soient un ensemble \( E\) de cardinal \( n\).
	\begin{enumerate}
		\item
		      si \( a\in E\) alors \( \Card(E\setminus\{ a \})=n-1\).
		\item		\label{ITEMooLAECooBLOpwZ}
		      Si \( a\not\in E\) alors \( \Card(E\cup\{ a \})=n+1\).
		\item		\label{ITEMooJTFGooKoBzgV}
		      \( \Card\big( \{ a \} \big)=1\).
		\item		\label{ITEMooTCMVooUTOqLw}<++>
		      \( \Card(\emptyset)=0\).
	\end{enumerate}
	%TODOooUTRHooZEZoUK. Prouver ça.
\end{lemma}


\begin{lemma}[\cite{MonCerveau}]	\label{LEMooNTVKooTfvHfA}
	Si \( E\) est un ensemble fini, alors son ensemble des parties \( \partP(E)\) est fini.
\end{lemma}

\begin{proof}
	Nous y allons par récurrence sur le cardinal\footnote{Voir \ref{NORMooPUOAooGnFjaX}.} de \( E\). Nous supposons que pour tout ensemble \( E'\) de cardinal \( n\), l'ensemble \( \partP(E')\) est fini. Soit maintenant \( E\) de cardinal \( n+1\).

	Soit \( a\in E\). La partie \( E'=E\setminus \{ a \}\) est de cardinal \( n\) par le lemme \ref{LEMooKNKEooYQeDdo}, donc \( \partP(E')\) est fini. En posant
	\begin{equation}
		\mS=\{ A\cup\{ a \}\tq A\in\partP(E') \},
	\end{equation}
	nous avons
	\begin{equation}
		\partP(E)=\partP(E')\cup \mS.
	\end{equation}
	De plus
	\begin{equation}
		\begin{aligned}
			\varphi\colon \mS & \to \partP(E')            \\
			B                 & \mapsto B\setminus\{ a \}
		\end{aligned}
	\end{equation}
	est une bijection. Donc \( \mS\) est fini et \( \partP(E)\) est fini comme union d'ensembles finis.
\end{proof}


\begin{proposition}[\cite{MonCerveau,BIBooFUKSooQARznK}]	\label{PROPooGJLYooTQtlSV}
	Tout ensemble fini partiellement ordonné\footnote{Définition \ref{DefooFLYOooRaGYRk}.} possède un élément maximal\footnote{C'est-à-dire un élément qui majore tous ceux avec qui il est comparable.}.

	Tout ensemble fini partiellement ordonné possède un élément minimal\footnote{C'est-à-dire un élément qui minore tous ceux avec qui il est comparable.}.
\end{proposition}

\begin{proof}
	Nous prouvons la contraposée, c'est-à-dire que si \( (E,\leq)\) n'a pas d'élément maximal, alors il est infini.
	\begin{subproof}
		\spitem[Axiome du choix, fonction \( s\)]
		%-----------------------------------------------------------
		Nous supposons que \( (E,\leq)\) n'admet pas d'élément maximal. Donc pour tout \( x\in E\) l'ensemble \( A_x=\{ y\in E\tq x<y \}\) est non vide. L'axiome du choix\footnote{Je suis certain qu'on peut réexprimer ça avec le lemme de Zorn; ce serait plus conforme à ce qui est fait dans le Frido, mais on ne va pas plus se casser la tête. Si \randomGender{le lecteur}{la lectrice} se sent capable de faire l'adaptation, qu'\randomGender{il}{elle} m'emvoie la preuve adaptée.} nous permet de choisir un élément dans chaque \( A_x\) et donc créer une application
		\begin{equation}
			s \colon E\to E
		\end{equation}
		telle que \( s(x)\in A_x\) pour tout \( x\).

		\spitem[La partie \( S\)]
		%-----------------------------------------------------------
		Soit \( x_0\in E\). Nous définissons la suite \( (x_n)_{n\in \eN}\) par récurrence\footnote{Théorème \ref{THOooEJPYooZFVnez}.} par \( x_{n+1}=s(x_n)\).

		\spitem[\( x_i< x_j\) si \( i<j\)]		\label{SPITEMooNUVJooDKClNE}
		%-----------------------------------------------------------
		Nous fixons \( i\in \eN\) et nous prouvons par récurrence sur \( j\) que \( x_i<x_j\) dès que \( i<j\). La définition \ref{DEFooAXZSooTEMjlV} de l'ordre sur \( \eN\) dit que si \( j>i\) alors il existe \( k\in \eN\) tel que \( j=i+k\). On y va par récurrence sur \( k\). D'abord nous avons
		\begin{equation}
			x_{i+1}=s(x_i)>x_i,
		\end{equation}
		donc \( x_{i+1}\neq x_i\). En ce qui concerne la récurrence, si \( x_{i+k}>x_i\), alors
		\begin{equation}
			x_{i+n+1}=s(x_{i+n})>x_{i+n}>x_i.
		\end{equation}
		Nous avons le résultat par transitivité de l'inégalité.

		\spitem[L'application \( \varphi\)]
		%-----------------------------------------------------------
		Nous posons \( X=\{ x_i \}_{i\in \eN}\), et nous posons
		\begin{equation}
			\begin{aligned}
				\varphi\colon E & \to E                            \\
				x               & \mapsto \begin{cases}
					                          s(x) & \text{si } s\in X \\
					                          x    & \text{sinon. }
				                          \end{cases}
			\end{aligned}
		\end{equation}

		\spitem[\( \varphi(E)\subset E\setminus\{ x_0 \}\)]
		%-----------------------------------------------------------
		D'abord notons que \( x_0\in X\), donc \( \varphi(x_0)=s(x_0)>x_0\), donc \( \varphi(x_0)\neq x_0\). Si \( y\in X\), alors il existe \( n\) tel que \( y=x_n\). Dans ce cas \( \varphi(y)=x_{n+1}>x_0\) par \ref{SPITEMooNUVJooDKClNE}. Si au contraire \( y\not\in X\) alors \( \varphi(y)=y\not\in X\) et en particulier \( y\neq x_0\).

		\spitem[\( E\setminus\{ x_0 \}\subset\varphi(E)\)]
		%-----------------------------------------------------------
		Si \( y\in X\setminus\{ x_0 \}\) alors \( y=x_n=s(x_{n-1})=\varphi(x_{n-1})\). Si au contraire \( y\) n'est pas dans \( X\) alors \( y=\varphi(y)\in \varphi(E)\).

		\spitem[\( E\setminus\{ x_0 \}=\varphi(E)\)]
		%-----------------------------------------------------------
		Ce sont les deux inclusions que nous venons de prouver.

		\spitem[Conclusion]
		%-----------------------------------------------------------
		L'application \(\varphi \colon E\to E\setminus\{ x_0 \}  \) est une bijection entre \( E\) et le sous-ensemble propre \( E\setminus\{ x_0 \}\). Donc \( E\) est infini.
	\end{subproof}

	En ce qui concerne l'énoncé à propos de l'existence d'un élément minimal, il suffit de prendre un élément maximal pour l'ensemble ordonné \( (E,\geq)\) (proposition \ref{PROPooNPLNooTqXVYz}.)
\end{proof}



%+++++++++++++++++++++++++++++++++++++++++++++++++++++++++++++++++++++++++++++++++++++++++++++++++++++++++++++++++++++++++++
\section{Les entiers relatifs}
\label{SECooNombresEntiers}
%+++++++++++++++++++++++++++++++++++++++++++++++++++++++++++++++++++++++++++++++++++++++++++++++++++++++++++++++++++++++++++

\begin{propositionDef}[\cite{RWWJooJdjxEK}]     \label{PROPooFIKUooVHlvTt}
	Soient \( a,b,a',b'\in \eN\). Nous disons que \( (a,b)\sim(a',b')\) si et seulement si
	\begin{equation}
		a+b'=b+a'
	\end{equation}
	\begin{enumerate}
		\item
		      \( \sim\) est une relation d'équivalence sur \( \eN^2\).
		\item       \label{ITEMooZQSHooSDfdvK}
		      Si \( (a,b)\sim (a',b')\) et \( (x,y)\sim (x',y')\) alors
		      \begin{equation}
			      (a+x,b+y)\sim(a'+x',b'+y').
		      \end{equation}
	\end{enumerate}
	L'ensemble des \defe{entiers}{entier} est
	\begin{equation}
		\eZ=(\eN\times \eN)/\sim,
	\end{equation}
	et nous notons \( \overline{ a,b }\in \eZ\) la classe de \( (a,b)\in \eN\times \eN\).
\end{propositionDef}

\begin{proof}
	En plusieurs points.
	\begin{subproof}
		\spitem[Symétrie]
		C'est la commutativité de la somme dans \( \eN\), proposition \ref{PROPooTLTSooGNMTmV}\ref{ITEMooIFFPooXfftfG}.
		\spitem[Réflexive]
		Immédiat.
		\spitem[Transitive]
		Nous supposons que \( (a,b)\sim(u,v)\) et que \( (u,v)\sim(x,y)\). Alors nous avons
		\begin{subequations}
			\begin{align}
				a+v & =u+b  \\
				u+y & =v+x.
			\end{align}
		\end{subequations}
		En additionnant membre à membre,
		\begin{equation}
			a+v+u+y=u+b+v+x.
		\end{equation}
		La commutativité nous permet de mettre \( u\) et \( v\) à droite dans chacun des deux membres. Ensuite la proposition \ref{PROPooTLTSooGNMTmV}\ref{ITEMooNUTHooJWWzGv} nous permet de simplifier par \( u+v\). Il reste \( a+y=b+x\), qui signifie \( (a,b)\sim(x,y)\).
		\spitem[Pour \ref{ITEMooZQSHooSDfdvK}]
		L'hypothèse donne les égalités
		\begin{subequations}
			\begin{align}
				a+b' & =b+a' \\
				x+y' & =y+x'
			\end{align}
		\end{subequations}
		En sommant, et en utilisant l'associativité,
		\begin{equation}
			(a+x)+(b'+y')=(b+y)+(a'+x').
		\end{equation}
		Cela signifie bien que \( (a+x,b+y)\sim(a'+x',b'+y')\).
	\end{subproof}
\end{proof}

\begin{lemma}	\label{LEMooEntiersEquivalentsZero}
	Soient \( a,b\in \eN\). Nous avons \( (a,b)\sim (0,0)\) si et seulement si \( a=b\).
\end{lemma}

\begin{proof}
	Dire que \( (a,b)\sim (0,0)\) est équivalent à dire que \( a+0=b+0\), ou encore que \( a=b\).
\end{proof}

\begin{propositionDef}[\cite{RWWJooJdjxEK}]	\label{PROPooDEFAdditionEntiersRelatifs}
	Soient \( a,b,x,y\in \eN\). L'application
	\begin{equation}
		\begin{aligned}
			f\colon \overline{ (a,b)\times \overline{ (x,y) } } & \to \eZ               \\
			\big( (a',b'),(x',y') \big)                         & \mapsto (a'+x',b'+y')
		\end{aligned}
	\end{equation}
	est constante.

	Nous nommons sa valeur \( \overline{ (a,b) }+\overline{ (x,y) }\).
\end{propositionDef}

\begin{proof}
	Cela est une conséquence de la proposition \ref{PROPooFIKUooVHlvTt}\ref{ITEMooZQSHooSDfdvK}.
\end{proof}

\begin{proposition}	\label{PROPooGroupeEntiersRelatifs}
	L'ensemble \( \eZ \), muni de l'addition telle que définie par la proposition \ref{PROPooDEFAdditionEntiersRelatifs}, est un groupe commutatif admettant pour neutre l'élément \( e = \overline{(0,0)} \).
\end{proposition}

\begin{proof}
	En plusieurs points.
	\begin{subproof}
		\spitem[Neutre]
		Pour tous entiers naturels \( a,\ b \), on a
		\begin{equation}
			\overline{ (a,b) }+\overline{ (0,0) }=\overline{ (a+0,b+0) }=\overline{ (a,b) }.
		\end{equation}
		De même \( e+\overline{ (a,b) }=\overline{ (a,b) }\) par commutativité de la somme dans \( \eN\).
		\spitem[Inverse]
		Il faut vérifier que \( \overline{ (b,a) }\) est l'inverse de \( \overline{ (a,b) }\).
		% TODOooPreuveOpposeRelatifs à faire.
		% TODOooDQGTooKXKrZY. Déjà dans ma liste.
		\spitem[Associativité]
		Calcul direct en utilisant l'associativité dans \( \eN\).
	\end{subproof}
\end{proof}

\begin{proposition}		\label{PROPooVMGYooUWsbSr}
	L'application
	\begin{equation}
		\begin{aligned}
			\iota\colon \eN & \to \eZ                    \\
			n               & \mapsto \overline{ (n,0) }
		\end{aligned}
	\end{equation}
	est un morphisme\footnote{Certes \( \eN\) n'est pas un groupe, donc le mot «morphisme» est un peu abusé, mais vous voyez ce que je veux dire.} injectif.
\end{proposition}

\begin{proof}
	Le fait que ce soit un morphisme est le calcul
	\begin{equation}
		\iota(a+b)=\overline{ (a+b,0) }=\overline{ (a,0) }+\overline{ (b,0) }=\iota(a)+\iota(b).
	\end{equation}

	Pour l'injectivité, supposons que \( \iota(a)=\iota(b)\). Alors \( \overline{ (a,0) }=\overline{ (b,0) }\), c'est-à-dire \( a+0=b+0\). Donc \( a=b\).
\end{proof}

%---------------------------------------------------------------------------------------------------------------------------
\subsection{Opposé}
\label{SUBooEntiersRelatifsOpposes}
%---------------------------------------------------------------------------------------------------------------------------

\begin{lemma}       \label{LEMooSABNooZZDIes}
	Tout élément de \( \eZ\) a un représentant de la forme \( (a,0)\) ou \( (0,b)\).
\end{lemma}

\begin{proof}
	Soient \( a,b\in \eN\). Si \( b\leq a\), alors nous avons
	\begin{equation}
		(a,b)\sim(a-b,0)
	\end{equation}
	où la différence est calculée dans \( \eN\) et a un sens parce que nous avons supposé \( b\leq a\). Si par contre \( a\leq b\) alors
	\begin{equation}
		(a,b)\sim(0,b-a).
	\end{equation}
	Puisque l'ordre sur \( \eN\) est total\footnote{Proposition \ref{PROPooGCCRooFBYrlo}.}, tous les cas sont couverts.
\end{proof}

\begin{lemmaDef}	\label{LEMooDEFRelatifOppose}
	Soit \( z\in \eZ\). L'application\footnote{Pour rappel, \( z\) est une classe d'équivalence dans \( \eN\times \eN\), c'est-à-dire une partie de \( \eN\times \eN\). Ça a un sens de prendre \( z\) comme ensemble sur lequel on définit une fonction.}
	\begin{equation}
		\begin{aligned}
			f\colon z & \to \eZ                    \\
			(a,b)     & \mapsto \overline{ (b,a) }
		\end{aligned}
	\end{equation}
	est constante.

	Nous nommons \( -z\) sa valeur.
\end{lemmaDef}

\begin{proof}
	Soient \( (a,b)\) et \( (x,y)\) dans \( z\). Nous avons successivement :
	\begin{itemize}
		\item
		      \( (a,b)\sim (x,y)\).
		\item
		      \( a+y=b+x\).
		\item
		      \( (b,a)\sim (y,x)\)
		\item
		      \( \overline{ (b,a) }=\overline{ (y,x) }\)
		\item
		      \( f(a,b)=f(x,y)\).
	\end{itemize}
	D'où la constance de \( f\).
\end{proof}

\begin{proposition}[\cite{MonCerveau}]	\label{PROPooJMETooHQGwnv}
	Soient \( a,b\in \eZ\). Nous avons
	\begin{enumerate}
		\item		\label{ITEMooWNYGooMopPcq}
		      \( (-a)b=-(ab)\)
	\end{enumerate}
\end{proposition}


\begin{lemma}	\label{LEMooEntiersRelatifsSigne}
	Nous avons
	\begin{enumerate}
		\item       \label{ITEMooSQFGooQPgIMu}
		      \( \eZ=\iota(\eN)\cup -\iota(\eN)\)
		\item       \label{ITEMooHQUQooJeqULl}
		      \( \iota(\eN)\cap -\iota(\eN)=\{ 0 \}\).
	\end{enumerate}
\end{lemma}

\begin{proof}
	Nous avons
	\begin{subequations}
		\begin{align}
			\iota(\eN)  & =\{ \overline{ (n,0) }\tq n\in \eN \}    \label{SUBEQooVJGVooCUxtvk} \\
			-\iota(\eN) & =\{ \overline{ (0,n) }\tq n\in \eN \}.   \label{SUBEQooAPGRooOCkYRr}
		\end{align}
	\end{subequations}
	Montrons à présent les deux points.
	\begin{subproof}
		\spitem[Pour \ref{ITEMooSQFGooQPgIMu}]
		Nous savons par le lemme \ref{LEMooSABNooZZDIes} que tous les éléments de \( \eZ\) sont de la forme \eqref{SUBEQooVJGVooCUxtvk} ou \eqref{SUBEQooAPGRooOCkYRr}.
		\spitem[Pour \ref{ITEMooHQUQooJeqULl}]
		Si \( z\in \iota(\eN)\cap -\iota(\eN)\), il existe \( n,m\in \eN\) tels que \( \overline{ (n,0) }=\overline{ (0,m) }\), ce qui signifie en particulier que \( (n,0)\sim(0,m)\) ou encore que \( n+m=0\). Le lemme \ref{LEMooQBHFooCuCusQ} dit alors que \( n=m=0\).

		Nous avons donc \( z= \overline{ (0,0) }=0\).
	\end{subproof}
\end{proof}


\begin{proposition}		\label{PROPooCHGRooRksGGO}
	L'application
	\begin{equation}
		\begin{aligned}
			\iota\colon \eN & \to \eZ                    \\
			n               & \mapsto \overline{ (n,0) }
		\end{aligned}
	\end{equation}
	est un morphisme injectif qui respecte l'ordre.
\end{proposition}

\begin{proof}
	Le fait que ce soit un morphisme injectif est la proposition \ref{PROPooVMGYooUWsbSr}.

	Il faut encore prouver que ça respecte l'ordre\quext{Faites-le en envoyez-moi la preuve.}
	%TODOooKQPQooAUkOlS il faut le faire.
\end{proof}


\begin{proposition}[\cite{MonCerveau}]	\label{PROPooORABooXRbVoz}
	Soit \( G\), un sous-groupe de \( \eZ\) non réduit à \( \{ 0 \}\). Alors
	\begin{enumerate}
		\item		\label{ITEMooUZBJooUeVmmz}
		      La partie \( \{ x>0\tq x\in G \}\) est non vide.
		\item	\label{ITEMooIXXFooOWkbPb}
		      En posant
		      \begin{equation}
			      p=\min\{ x>0\in G \},
		      \end{equation}
		      nous avons \( G=p\eZ\).
	\end{enumerate}
\end{proposition}

\begin{proof}
	Vu que \( G\) est un sous-groupe de \( (\eZ,+)\), si \( x,y\in G\) nous avons \( x+y\in G\), \( -x\in G\), etc.
	\begin{subproof}
		\spitem[Pour \ref{ITEMooUZBJooUeVmmz}]
		%-----------------------------------------------------------
		Soit \( x\neq 0\) dans \( G\). Si \( x>0\) c'est bon. Sinon \( -x\in G\) et \( -x>0\).
		\spitem[Pour \ref{ITEMooIXXFooOWkbPb}]
		%-----------------------------------------------------------
		Nous posons
		\begin{equation}
			p=\min\{ x\in G\tq x>0 \}.
		\end{equation}
		Vu que \( p\in G\) nous avons \( p\eZ\subset G\). Nous devons montrer l'inclusion inverse.

		Soit \( y\in G\). Par la division euclidienne \ref{THOooKDJVooRIJRHP}, il existe \( q\in \eZ\) et \( 0\leq r<p\) tels que \( y=qp+r\). Vu que \( y\) et \( qp\) sont dans \( G\) nous avons \( r=y-qp\in G\). Par minimalité de \( p\), nous devons avoir \( r=0\), et donc \( y=qp\in p\eZ\).
	\end{subproof}
\end{proof}

%---------------------------------------------------------------------------------------------------------------------------
\subsection{Ordre sur \( \eZ\)}
\label{SUBooEntiersRelatifsOrdre}
%---------------------------------------------------------------------------------------------------------------------------

\begin{normaltext}	\label{NORMooNotationInclusionEntiers}
	Si \( z\in \eZ\), nous disons que \( z\in \eN\) lorsque \( z\in \iota(\eN)\). C'est un abus de notation qu'il est difficile de ne pas faire. Le même abus sera fait pour passer de \( \eZ\) à \( \eQ\) puis \( \eR\) puis \( \eC\).
\end{normaltext}

% TODOooKCJMooIHGrZN donner une notion de positivité sur Z, et modifier la définition PROPooMYYDooOABOdB.
% Fusionner
% - DEFooYMRHooCyEKtr
% - DEFooRelatifsPositifs
% - PROPooMYYDooOABOdB
% - PROPooMYYDooOABOdB
% ceci est déjà dans ma liste

\begin{definition}[Positivité et ordre dans \( \eZ\)\cite{MonCerveau}]	\label{DEFooYMRHooCyEKtr}
	La positivité\footnote{Définition \ref{DEFooZRMFooCtzMov}.} que nous considérons sur \( \eZ\) est \( \eN\setminus\{ 0 \}\). Nous considérons alors la relation d'ordre\footnote{De la proposition \ref{PROPooKLOPooBgQqhM}.} \( \leq\) donnée par \( a\leq b\) si et seulement si \( b-a\in \eN\).

	Nous disons que
	\begin{enumerate}
		\item
		      L'élément \( z\in \eZ\) est \defe{positif}{entier relatif!positif} si \( z\geq 0\).
		\item
		      L'élément \( z\in \eZ\) est \defe{strictement positif}{entier relatif!strictement positif} si \( z\geq 0\) et \( z\neq 0\). Nous notons \( z>0\).
		\item
		      L'élément \( z\in \eZ\) est \defe{négatif}{entier relatif!négatif} si \( z\leq 0\).
		\item
		      L'élément \( z\in \eZ\) est \defe{strictement négatif}{entier relatif!strictement négatif} si \( z\leq 0\) et \( z\neq 0\). Nous notons \( z<0\).
	\end{enumerate}
\end{definition}

\begin{definition}[Éléments positifs de \( \eZ \)]	\label{DEFooRelatifsPositifs}
	On dit qu'un élément \( z \in \eZ \) est \defe{positif}{entiers relatifs!positifs} si \( z \in \eN \). On note dans ce cas \( z \geq 0\).

	On dira qu'il est \defe{négatif}{entiers relatifs!négatifs} quand \( -z \in \eN \). On note ainsi \( z \leq 0 \).

	On dit qu'un entier relatif \( z \) est \defe{strictement positif}{} s'il est positif et non nul; on note \( z > 0 \). Il est \defe{strictement négatif} lorsqu'il est négatif et non nul, et cela est noté \( z < 0 \).
\end{definition}

\begin{propositionDef}[Relation d'ordre \cite{RWWJooJdjxEK}]       \label{PROPooMYYDooOABOdB}
	La paire\footnote{Voir la définition \ref{DEFooYMRHooCyEKtr} pour \( \leq\).} \( (\eZ, \leq)\) est un ensemble totalement ordonné.
\end{propositionDef}

\begin{normaltext}
	%\label{NORMooOrdreEntiersRelatifsExtension}
	% Ce label me semble dangereux parce que si on le cite dans une démonstration entre ici et LEMooKAXFooIPyzJC, le détecteur de référence vers le futur
	% ne préviendra pas du problème.
	Une version de la proposition \ref{PROPooMYYDooOABOdB} dans \( \eR\) sera le lemme \ref{LEMooKAXFooIPyzJC}.
\end{normaltext}


\begin{proposition}[\cite{MonCerveau}]	\label{PROPooADGIooTioqdu}
	L'anneau \( \eZ\) est totalement ordonné\footnote{Définition \ref{DEFooWACWooDWvXKJ}.}
	%TODOooBRCUooRjkxmG. Prouver ça.
	% Pour ça, il faudra encore définir la multiplication dans Z comme extension de celle dans N.
\end{proposition}

\begin{lemma}       \label{LEMooSVDDooWsyxNP}
	Soient \( a>0\) et \( b \geq 1\) (resp. \( b > 1 \) dans \( \eZ\). Nous avons
	\begin{equation}
		ab \geq a \qquad \text{(resp. }ab>a\text{)}.
	\end{equation}
\end{lemma}

\begin{proof}
	On la fait seulement pour l'inégalité stricte: il suffit de recopier pour l'inégalité large. Par récurrence sur \( a \) qui se trouve être dans \( \eN \). Lorsque \( a = 1 \), on lit simplement \( 1b > 1 \) qui est correct.

	Si la propriété est vraie pour un certain \( a \), alors
	\begin{align}
		(a+1) b & = ab + b                          &  & \text{(distributivité)} \\
		        & > a + b\text{(hyp. de réc.)}                                   \\
		        & > a + 1\text{(hyp. sur \( b \)).}
	\end{align}
	Et le principe de récurrence permet de conclure.
\end{proof}

\begin{normaltext}	\label{NORMooLienRelatifsAnnTotOrdonne}
	Cette proposition est en partie similaire à \ref{PROPooELXCooCYzEVD} qui parle d'anneau totalement ordonné.
\end{normaltext}

\begin{proposition}[\cite{MonCerveau}]	\label{PROPooYFUBooJUZgwH}
	Quelques propriétés de l'ordre dans \( \eZ\). Soit \( a, b \in \eZ\).
	\begin{enumerate}
		\item		\label{ITEMooULKWooLmKBse}
		      Si \( a,b\geq 0\), alors \( ab\geq 0\).
		\item	\label{ITEMooSAJPooEDSgXJ}
		      Si \( a,b>0\), alors \( ab>0\).
		\item		\label{ITEMooINGNooOgSGiF}
		      Si \( a,b\leq 0\) alors \( ab\geq 0\).
		\item	\label{ITEMooOGFXooDvLKAE}
		      Si \( a\geq 0\) et \( b\leq 0\) alors \( ab\leq 0\).
		\item	\label{ITEMooRelatifsOrdreOppose}
		      Si \( a>0\), alors \( -a<0\).
		\item	\label{ITEMooRelatifsOrdreCarrePositif}
		      Nous avons \( a^2\geq 0\) et \( a^2=0\) si et seulement si \( a=0\).
		\item		\label{ITEMooBUHIooLJOSLQ}
		      Si \( ab\geq 0\) et si \( a\geq 0\) alors \( b\geq 0\).
		\item		\label{ITEMooGDFPooBbpBgn}
		      Si \( ab>0\) alors soit \( a,b>0\), soit \( a,b<0\).
	\end{enumerate}
	%TODOooEDWPooRDujTs. Prouver ça. J'en mets deux parce que je parie que ce sera fait.
\end{proposition}

\begin{lemma}       \label{LEMooMYEIooNFwNVI}
	Toute partie bornée de \( \eZ\) possède un plus grand élément.
	%TODOooEDWPooRDujTs. Prouver ça. J'en mets deux parce que je parie que ce sera fait.
\end{lemma}

\begin{proposition}	\label{PROPooRelatifsValAbs}
	Si \( z \in \eZ \), alors :
	\begin{enumerate}
		\item
		      \( |z| \) est bien définie et \( |z| \in \eN \);
		\item
		      nous avons \( | ab |=| a || b | \).
	\end{enumerate}
\end{proposition}

\begin{proof}
	C'est le fait que \( \eZ\) est un anneau totalement ordonné\footnote{Proposition \ref{PROPooADGIooTioqdu}.}: donc la proposition-définition \ref{DEFooJXKVooErANPh} s'applique, y compris l'équation \eqref{EQooORNAooBYFeFk}.
\end{proof}


\begin{definition}	\label{DEFooRelatifsDiviseurParite}
	Soit \( a, b \in \eZ \). On étend les notions de divisibilité et de parité de la définition \ref{DEFooDivisibleNombrePairImpair} en disant que:
	\begin{enumerate}
		\item
		      \( a \) est divisible par \( b \) si \( |a| \) est divisible par \( |b| \);
		\item
		      \( a \) est pair si \( |a| \) est pair;
		\item
		      \( b \) est impair si \( |b| \) est impair.
	\end{enumerate}
\end{definition}

\begin{proposition} \label{PROPooYJBMooZrzkNX}
	Soit \( a,b\in \eZ\) non nuls, tels que \( a\) soit divisible par \( b\). Alors \( | a |\geq | b |\).
\end{proposition}

\begin{proof}
	\( a\) est divisible \( b\) dans \( \eZ \) revient à dire que \( |a| \) est divisible \( |b| \) dans \( \eN \), donc par la définition \ref{DEFooDivisibleNombrePairImpair}, il existe \( q \in \eN \) tel que \( |a| = q\, |b| \). Comme \( a \) et \( b \) sont non nuls, \( |a| \) et \( |b| \) sont aussi non nuls, et donc \( q \neq 0 \): en effet, sinon, nous aurions \( q\, |b| = 0 = |a| \) à cause de la proposition \ref{PROPooGHDOooFYRmon} point \ref{ITEMooNBYKooXnGRrf}. Ainsi, \( q \geq 1 \), et grâce au lemme \ref{LEMooSVDDooWsyxNP}, on a \(|a| = q\, |b| \geq |b| \).
\end{proof}

\begin{lemma}[\cite{MonCerveau}]	\label{LEMooCMOEooSoeclk}
	Soit \( n\in \eZ\).
	\begin{enumerate}
		\item
		      Si \( n\) est impair\footnote{Définition \ref{DEFooRelatifsDiviseurParite}.}, alors \( n-1\) et \( n+1\) sont pairs.
		\item
		      Si \( n\) est pair, alors \( n-1\) et \( n+1\) sont impairs.
	\end{enumerate}
\end{lemma}

\begin{proof}
	Selon que \( n \geq 0 \) ou que \( n \leq 0 \), on effectue la division euclidienne de \( n \) ou de \( -n \) par 2.
	\begin{subproof}
		\spitem[Cas \( n \geq 0 \)]
		%--------------------------------------------------------
		Il existe \( q \in \eN \) et \( r \in \{ 0, 1\} \) tel que \( n = 2q+r \).
		\begin{itemize}
			\item
			      Quand \( n \) est pair, on a \( r = 0 \). Alors \( n+1 = 2q + 1 \) et \( n - 1 = 2(q-1) + 2 - 1 =  2(q-1) + 1 \) sont impairs.
			\item
			      Quand \( n \) est impair, on a \( r = 1 \). Alors \( n+1 = 2q + 2 = 2(q+1) \) et \( n - 1 = 2q + 1 - 1 =  2q \) sont pairs.
		\end{itemize}
		\spitem[Cas \( n \leq 0 \)]
		%--------------------------------------------------------
		Il existe \( q \in \eN \) et \( r \in \{ 0, 1\} \) tel que \( -n = 2q+r \).
		\begin{itemize}
			\item
			      Quand \( n \) est pair, on a \( r = 0 \). Alors \( -n+1 = 2q + 1 \) est impair, et \( n-1 = -(-n+1) \) l'est aussi. De même, \( -n - 1 = 2(q-1) + 2 - 1 =  2(q-1) + 1 \), et donc \( n+1 = -(-n-1) \), sont impairs.
			\item
			      Quand \( n \) est impair, on a \( r = 1 \). Alors \( -n+1 = 2q + 2 = 2(q+1) \) est pair, et \( n-1 = -(-n+1) \) l'est aussi. De même, \( -n - 1 = 2q + 1 - 1 =  2q \), et donc \( n+1 = -(-n-1) \), sont pairs.
		\end{itemize}
	\end{subproof}
\end{proof}

\begin{lemma}       \label{LEMooJNXIooBmdOVi}
	L'ensemble \( \eZ\) est infini dénombrable.
\end{lemma}

\begin{proof}
	Nous considérons l'application
	\begin{equation}
		\begin{aligned}
			f\colon \eN & \to \eZ                                         \\
			n           & \mapsto \begin{cases}
				                      n/2      & \text{si \( n\) est pair}    \\
				                      -(n+1)/2 & \text{si \( n\) est impair,}
			                      \end{cases}
		\end{aligned}
	\end{equation}
	et nous montrons qu'elle est bijective.
	\begin{subproof}
		\spitem[\( f\) est injective]
		%-----------------------------------------------------------
		Soient \( x,y\in \eZ\) tels que \( f(x)=f(y)\). Il y a trois possibilités : \( f(x)<0\), \( f(x)=0\) et \( f(x)>0\).
		\begin{subproof}
			\spitem[Si \( f(x)=0\)]
			%-----------------------------------------------------------
			Nous supposons que \( f(x)=0\). Si \( x \) est pair, c'est que \( x/2=0\) et donc que \( x=0\). Si \( x\) est impair, c'est que \( x+1=0\) et donc que \( x=-1\) qui n'est pas possible parce que \( x\in \eN\). Bref, la seule possibilité pour \( f(x)=0\) est \( x=0\).

			\spitem[Si \( f(x)<0\)]
			%-----------------------------------------------------------
			Alors \( x\) est pair et \( x=2f(x)\). Même raisonnement pour \( y\) : \( y=2f(y)=2f(x)=x\).

			\spitem[Si \( f(x)<0\)]
			%-----------------------------------------------------------
			Alors \( x\) est impair et
			\begin{equation}
				f(x)=-\frac{ x+1 }{ 2 }=f(y)=-\frac{ y+1 }{ 2 },
			\end{equation}
			et donc \( x=y\).
		\end{subproof}

		\spitem[\( f\) est surjective]
		%-----------------------------------------------------------
		Soit \( n\in \eZ\). Si \( n>0\), alors \( n=f(2n)\). Si \( n<0\) alors \( n=f(-2n-1)\). Notez que \( -2n-1>0\), donc c'est valide.
	\end{subproof}
\end{proof}

\begin{lemma}[\cite{MonCerveau}]	\label{LEMooILYLooTDRtYj}
	Soit un anneau\footnote{Définition \ref{DefHXJUooKoovob}.} \( A\). Si une application \(f \colon \eZ\to A  \) vérifie
	\begin{enumerate}
		\item
		      \( f(x+y)=f(x)+f(y)\) pour tout \( x,y\in \eZ\),
		\item
		      \( f(1)=1\),
	\end{enumerate}
	alors
	\begin{enumerate}
		\item		\label{ITEMooUQKWooByNDTK}
		      \( f(0)=0\)
		\item	\label{ITEMooYOEPooKOQQPH}
		      \( f(-x)=-f(x)\)
		\item	\label{ITEMooKTSDooExAJHd}
		      \( f\) est un morphisme d'anneaux, c'est-à-dire que
		      \begin{equation}
			      f(xy)=f(x)f(y).
		      \end{equation}
	\end{enumerate}
\end{lemma}

\begin{proof}
	En deux parties.
	\begin{subproof}
		\spitem[Pour \ref{ITEMooUQKWooByNDTK}]
		%-----------------------------------------------------------
		Prenons un \( x\) quelconque. Nous avons \( f(x)= f(x+0)=f(x)+f(0)\). Donc \( f(0)=0\).

		\spitem[Pour \ref{ITEMooYOEPooKOQQPH}]
		%-----------------------------------------------------------
		L'inverse pour la somme de \( f(-x)\) est \( -f(x)\). En effet
		\begin{equation}
			f(-x)+f(x)=f(-x+x)=f(0)=0.
		\end{equation}
		\spitem[Pour \ref{ITEMooKTSDooExAJHd}, si \( y>0\)]
		%-----------------------------------------------------------
		Nous prouvons que \( f(xy)=f(x)f(y)\) par récurrence sur \( y\). D'abord pour \( y=1\) nous avons bien \( f(xy)=f(x)f(y)\) parce que \( f(1)=1\). Ensuite nous avons :
		\begin{subequations}
			\begin{align}
				f\big( x(y+1) \big) & =f(xy+x)                                         \\
				                    & =f(xy)+f(x)                & \text{morph pour +} \\
				                    & =f(x)f(y)+f(x)             & \text{hyp. rec.}    \\
				                    & =f(x)\big( f(y)+1 \big)                          \\
				                    & =f(x)\big( f(y)+f(1) \big)                       \\
				                    & =f(x)f(y+1)
			\end{align}
		\end{subequations}
		\spitem[Si \( y<0\)]
		%-----------------------------------------------------------
		Nous posons \( y'=-y\). Vu que \( y'>0\) nous avons
		\begin{equation}
			f(xy)=f(-xy')=-f(xy')=-f(x)f(y')=f(x)(-1)f(y')=f(x)f(-y')=f(x)f(y).
		\end{equation}
	\end{subproof}
\end{proof}
